\subsection{Geometric classes and the cyclotomic centralizer}

We retain the tangent at $O$ with parameter $z/x$ and tangent one.
All inertia groups in this section are their images in the geometric
pro-$3$ group $N$.  We write $\kappa$ for continuous pro-$3$ Kummer
classes.  A cuspidal section means a section contained in the
decomposition group of a lifted cusp, with geometric conjugacy taken
in the group actually displayed.  The manuscript comparison is
\S51.16, folios 733v--736v, scans 172--178, formulas (371)--(379).

\begin{statement}{B.1}{Proved}{The three abelian cusp packets, with their framed origins}
Put $J=J_{E,D}$ and use the frame coordinates of LM7 (E24.4).
Let $z_P$ be the framed line bundle $\mathcal O(P-O)$ with frames
$(x/z)^{-1}$ at $P$, $z/x$ at $O$, and the canonical section at the
other cusps.  Define $z_Q$ in the same way from $Q_c-O$, using $y/z$
at $Q_c$.  Put $\alpha_P=\kappa(z_P)$, $\alpha_Q=\kappa(z_Q)$ and
$\alpha_O=0$.  The three packets in $H^1(G_{\mathbb Q},V)$ are
\[
 \mathscr P_d^V=\alpha_d+e_d\,\mathcal R_{\mathbb P}(\mathbb Z_3),
 \qquad d=O,P,Q_c.                                      \tag{B.1}
\]
Here $e_d\,\mathcal R_{\mathbb P}(\mathbb Z_3)$ denotes the image of
the continuous homomorphism on cohomology induced by the
$G_{\mathbb Q}$-equivariant inclusion
$\mathbb Z_3(1)\to V$, $1\mapsto e_d$.

The torus-point relation specifying the second framed origin is
\[
 z_P+z_Q=
 (3^{-1},3^{-1},-\lambda^{-1},-\lambda^{-1},\zeta\lambda^{-2}).
                                                               \tag{B.2}
\]
In the presentation (E24.7), eliminate the first $\lambda$-generator
and retain $z'$.  Write $g_2,g_3$ for the remaining two
$\lambda$-generators and $\tau_3$ for $\zeta$ in the third quadratic
factor.  After pro-$3$ completion the origins have toric coordinates
zero, divisible-cohomology coordinate zero, and remaining coordinates
\[
 \alpha_P=\kappa(z'),\qquad
 \alpha_Q=\kappa(2z'+g_2+g_3-\tau_3).                   \tag{B.3}
\]
The inertia maps in the original frame coordinates are
\[
 \begin{split}
 i_O(u)&=(u^{-1},u^{-1},u^{-1},u^{-1},u^{-1}),\\
 i_P(u)&=(u,1,1,1,1),\qquad i_Q(u)=(1,u,1,1,1).
 \end{split}                                               \tag{B.4}
\]
Their images under the two coordinate-unit characters are respectively
$(-\kappa u,-\kappa u)$, $(\kappa u,0)$ and $(0,\kappa u)$.
The three sets in (B.1) are pairwise disjoint, and each displayed
parameter map is injective.
\end{statement}
\begin{proof}
The tangential Albanese class at $d$, relative to the tangent at $O$,
is represented in the generalized Picard group by
$\mathcal O(d-O)$ with its local parameter frames.  This description
can be obtained directly from the divisor construction: replace $d$
and $O$ by points with local parameters $s_d$ and $s_O$, take the
line bundle of their difference, and divide its local frames by the
leading powers $s_d^{-1}$ and $s_O$.  The resulting fibers at the
reduced boundary are the displayed frames.  Applying multiplication
by $3^n$ identifies the finite abelian path torsor with the torsor of
$3^n$-division points of this framed line bundle.  The inverse limit
therefore gives its Kummer class in $T_3J=V$.
Changing the nonzero tangent at $d$ multiplies its local root system
by the root system of a scalar $u$ and changes the abelian cocycle by
$e_d\kappa(u)$.  The same description applies to an arbitrary
continuous section of the cusp decomposition extension, because that
extension is $\mathbb Z_3(1)\rtimes G_{\mathbb Q}$ and its section
classes are $H^1(G_{\mathbb Q},\mathbb Z_3(1))$.

We verify the origin calculation.  The rational function
\[
 h=\frac{x+y+z}{x+y}
\]
has divisor $P+Q_c-2O$.  Its leading values are $-1$ after division
by $x/z$ at $P$, $-1$ after division by $y/z$ at $Q_c$, and $-3$
after multiplication by $(z/x)^2$ at $O$.  On the representatives
$[0:-\zeta:1]$, $[-\zeta:0:1]$ and $[1:-\zeta:0]$ its values,
normalized by $-3$, are
$-\lambda^{-1},-\lambda^{-1},\zeta\lambda^{-2}$.
These are exactly (B.2).  The expansions follow, for example, from
$y/z=-(1+(x/z)^3)^{1/3}$ near $P$.  The relation
$3=-\zeta^2\lambda^2$ verifies the three quadratic values without
choosing complex approximations.

For comparison, multiplication by
$(x+z)/(x+y)$ trivializes $\mathcal O(3Q_c-3O)$ and gives
\[
 3z_Q=(3^{-1},-3^{-2},-\lambda^{-2},-\lambda^{-1},
                                  \zeta\lambda^{-3}).       \tag{B.5}
\]
Together with (E24.6), this checks the cube of (B.2) in all five
coordinates.  The program \texttt{code/bc\_checks.py} verifies all
these identities over the exact field $\mathbb Q[\zeta]$.

Signs disappear in pro-$3$ Kummer theory: $-1$ has compatible
rational $3^n$-th roots.  Relation (E24.7) then reads
$3z'=-g_1-2g_2-3g_3+2\tau_3$.  The quadratic coordinates in (B.2)
are $-g_1-g_2-2g_3+\tau_3$.  Subtracting $z'$ and eliminating $g_1$
gives $2z'+g_2+g_3-\tau_3$, proving (B.3).
The last factor of (E24.1) is zero because these are Kummer classes
of actual rational generalized-Jacobian points; we use the splitting
there with its rational-point completion as kernel.

Normalizing the frame at $O$ to one gives (B.4).  Applying (E24.4)
gives the three asserted toric images; the quadratic norm of a
rational scalar is its square.  These images prove injectivity of
each parameter map.  Filling all cusps gives the continuous
$G_{\mathbb Q}$-equivariant quotient $V\twoheadrightarrow T_3E$.
It sends the packets to the point classes of $O,P,Q_c$, respectively.
Their pairwise differences have exact order three.  Reduction at two
embeds rational odd torsion into $E(\mathbb F_2)$, which has order
three, as checked by the program.  Consequently none of these
nonzero differences is three-divisible in $E(\mathbb Q)$, and its
modulo-three Kummer class is nonzero.  This proves disjointness.
\end{proof}

\begin{statement}{B.7}{Proved}{Every nontrivial tame inertia in a full section is cuspidal}
Let $q\ne3$ be a rational prime, put $k=\mathbb Q_q$, and
use the actual pro-three fundamental group $N$ of the marked
Fermat cubic with the integral unit tangent at $O$ as origin.
For a continuous section $s:G_k\to N\rtimes G_k$, its
restriction to inertia is either trivial or has image in one
cuspidal inertia subgroup.  In the latter case the whole
section is contained in that cusp's decomposition group.
The cusp is $k$-rational and the lifted cusp is unique.
Moreover, a full section is unramified at $q$ exactly when
its image in $V$ is unramified there.
\end{statement}
\begin{proof}
The marked curve and the tangent have good reduction at $q$.
Thus the original based action of inertia on $N$ is trivial,
by prime-to-$q$ specialization.  Its cocycle restriction is
a homomorphism.  Wild inertia, being pro-$q$, maps to one;
the only tame factor which can map to $N$ is $\mathbb Z_3(1)$.
Let $t=s(\tau_q)$ denote the geometric value on a generator
of that factor, and assume $t\ne1$.  If
$F=s(\phi_q)$ is the full Frobenius lift, the section relation
is the exact group equality
\[
                         FtF^{-1}=t^q.                      \tag{B.7.1}
\]

Take a characteristic open normal subgroup $M$ of $N$ and put
$H=M\overline{\langle t\rangle}$.  Such $M$ form a cofinal
family, since $N$ is finitely generated.  Equation (B.7.1)
and the characteristic property of $M$ show that $H$ is
stable under $s(G_k)$.  The open arithmetic subgroup
$H\,s(G_k)$ is therefore the subgroup of the pro-three
arithmetic fundamental group defining a finite etale cover
$Y_H\to U_k$.  Its geometric group is $H$, by the exact
finite-three-group cover comparison proved in B.4.

We justify the reduction property needed for weights.
The subgroup $H\,s(G_k)$ contains the original arithmetic
inertia: its geometric cocycle values are powers of $t$,
which lie in $H$.  Original inertia centralizes $N$, so it
acts trivially on every coset in $N/H$.
The cover consequently descends to the smooth marked
integral model by prime-to-$q$ specialization.  Its
compactification is smooth and proper over the valuation
ring.  Locally at a boundary point this last assertion is
the tame normalization $z^e=x$ after an unramified extension,
with $e$ prime to $q$; away from the boundary it is etale.
These charts are exactly the tame-cover specialization
comparison, also obtained from the logarithmic comparison
in Stix, Theorem 3.4.6 \cite{BC_Stix2002}.

Let $X_H$ be this compactification.  The boundary sequence
(B.8), with $\mathbb Z_3$ coefficients, is equivariant for
the actual arithmetic Frobenius $F$ of this descended cover.
The image of $t$ in $T_3J_{X_H}$ satisfies $Fv=qv$ by
(B.7.1).  All Frobenius eigenvalues on this Tate module
have complex absolute value $\sqrt q$, by Deligne,
Theorem (1.6) \cite{BC_Deligne1974}, applied in degree one
to the smooth proper special fiber.  The Tate module of
its Jacobian has these same eigenvalues for arithmetic
Frobenius, by the Abel--Jacobi comparison and the dual
Weil pairing.  Thus $v=0$.
The image of $t$ in $H^{\rm ab}$ is in the boundary submodule.

The cyclic finite three-group $H/M$ is generated by the image
of $t$.  Its homomorphism from $H^{\rm ab}$ consequently maps
the boundary submodule onto all of $H/M$.  At least one
individual cusp inertia maps onto this cyclic group: otherwise
every such image would be in its unique maximal subgroup.
It follows that $t$ fixes a cusp in the boundary set of the
normal cover associated to $M$.
The inverse images of these finite fixed-cusp sets in
$\coprod_d N/I_d$ are nonempty closed subsets with the
finite-intersection property.  Compactness gives a lifted
cusp fixed by $t$ at every level, so $t$ belongs to its
actual inertia group $I_{\widetilde d}$.

Distinct cusp inertia groups have trivial intersection by
the boundary and cover argument in B.4.  Hence this lifted
cusp is unique.  Equation (B.7.1) shows that $F$ normalizes
$I_{\widetilde d}$, because its image has the nontrivial
element $t^q$ in common with that group.  Its normalizer is
the cusp decomposition group, again by the explicit
stabilizer argument in B.4.  The inertia image and $F$
topologically generate the image of the local section;
therefore that whole image lies in the decomposition group.
Its projection is all $G_k$, forcing $d\in D(k)$.

Finally the inclusion $I_{\widetilde d}\to V$ sends its
generator to the primitive nonzero boundary vector $e_d$.
It is injective.  Thus $t\ne1$ has nonzero abelian image,
while a trivial full inertia image certainly has trivial
abelian image.  This proves the last assertion.
\end{proof}

\begin{statement}{B.8}{Proved}{The exact local image of ramified full sections in class two}
With $k=\mathbb Q_q$, $q\ne3$, let $A=N/\gamma_3N$.
For each $d\in D(k)$ choose the coordinate unit tangent
used in C.10, its actual section $s_{d,1}$, and a peripheral
generator $c_d$ based at $O$.  The ramified part of the image
of full sections in $\operatorname{Sec}_k(A)$ is precisely
the disjoint union
\[
 \coprod_{d\in D(k)}
 \left\{
  \left[g\longmapsto
    c_d^{\kappa_q(u)(g)}s_{d,1}(g)\right]:
  u\in k^{\times\wedge_3},\ v_q(u)\ne0
 \right\}.                                                   \tag{B.8.1}
\]
Here the expression is multiplication in the arithmetic
group, followed by its actual class-two quotient.
Its parameter space is explicitly
\[
 k^{\times\wedge_3}
   =\mathbb Z_3\kappa(q)\oplus\mu_{3^{v_3(q-1)}}(k).
                                                               \tag{B.8.2}
\]
There are nine packets for $q\equiv1\pmod3$ and three for
$q\equiv2\pmod3$.  In the central cuspidal lattice
$I/[N,I]=\mathbb Z_3(1)[D]$ the tame vector of a packet is
$v_q(u)c_d$, with exactly this one cusp coordinate nonzero.
In particular the mixed-inertia counterexamples in A.5 and
A.6 cannot lift to full $N$.
\end{statement}
\begin{proof}
The decomposition extension for a rational cusp is
$\mathbb Z_3(1)\rtimes G_k$ using its unit tangent.
Its sections modulo inertia conjugacy are the Kummer group
(B.8.2).  The decomposition inclusion is a continuous
homomorphism of actual arithmetic groups; after quotienting
the geometric group it gives exactly (B.8.1).
The Kummer cocycle equation and
$s_{d,1}(g)c_ds_{d,1}(g)^{-1}=c_d^{\chi_3(g)}$
verify directly that the displayed product is a section.
The unit-tangent section is unramified, and the tame
Kummer exponent is $v_q(u)$, giving the stated ramification
condition.  Every displayed class has its full cuspidal
lift.  Conversely B.7 puts every ramified full section in
one of these decomposition extensions.

To check that the displayed parametrization remains
injective in class two, first pass to a finite unramified
extension splitting all cusps.  The boundary inclusion
$e_d:\mathbb Z_3(1)\to W$ is then split, and
$H^1(k,W)\to H^1(k,V)$ is injective because
$T_3E^{G_k}=0$.  Restriction on the Kummer group is
injective by the same invariant argument as C.10.
Thus two parameters at the same cusp cannot coincide.
Different cusps map in the proper elliptic quotient to
the Kummer classes of different three-torsion points.
Reduction at $q\ne3$ injects three-primary torsion into
the finite group $E(\mathbb F_q)$; its nonzero elements
are not divisible by every power of three.  The formal
group is pro-$q$, on which multiplication by three is
invertible.  Hence a nonzero difference of two cusp
points has nonzero integral local Kummer class, proving
disjointness.

The decomposition of $k^\times$ into a uniformizer,
the roots of unity of order $q-1$, and principal units
gives (B.8.2), since the principal units are pro-$q$.
The two possible cusp counts follow from the actual
field of the six nonrational cusps, $\mathbb Q(\zeta_3)$.
Projection to $I/[N,I]$ sends the peripheral generator
to its indicated basis element, proving the tame-vector
formula.  The vectors in A.5 and A.6 involve
$c_O+c_P$ and therefore do not have this form.
\end{proof}

\begin{statement}{B.9}{Proved}{All inherited class-two cusp origins}
Use exactly the diagonal coordinates of LM7 E.29 for $H$,
based at the lifted cusp $P$ over $t=0$, and put
\[
       f(g)=\frac{\chi_3(g)^2-1}{24}\in\mathbb Z_3(2).
                                                               \tag{B.9.1}
\]
The three rational cusp packets in $\operatorname{Sec}(H)$
have the following continuous cocycle representatives:
\[
 \begin{array}{c|ccc}
 &u(g)&v(g)&w(g)\\ \hline
 P&\kappa(a)(g)&0&0\\
 Q_c&0&\kappa(a)(g)&f(g)\\
 O&-\kappa(a)(g)&-\kappa(a)(g)&f(g)/2
 \end{array}
 \qquad a\in\mathbb Q^{\times\wedge_3}.                    \tag{B.9.2}
\]
Each row is injectively parametrized and the rows are
pairwise disjoint.  Above inherited abelian class $(0,0)$
the three unit tangents give, respectively, the central
labels $0,2,1\in\mathbb Z/3$ relative to
$g\mapsto(\chi_3(g)^2-1)/3$.
The actual arithmetic map from the intrinsic group based
at $O$ to this $H$-extension is
\[
 (n,g)\longmapsto
    \bigl(q_O(n)(0,0,f(g)/2),g\bigr),                       \tag{B.9.3}
\]
after the geometric path choices in the proof.
Thus (B.9.2) retains the change of section origin between
the $O$ and $P$ base points.
\end{statement}
\begin{proof}
Write $t=-u_0^3$, $1-t=-v_0^3$ for the cubic cover of the
tripod, with $u_0=x/z$, $v_0=y/z$.
The standard path from the unit tangent at zero to the
unit tangent at one lifts from $P$ to $Q_c$: along the
real interval its cube roots are
$u_0=-t^{1/3}$ and $v_0=-(1-t)^{1/3}$.
It reaches exactly the displayed rational cusps.
At either end the lifted local parameter has tangent
minus one relative to $u_0$ or $v_0$.
Changing that sign has zero pro-three Kummer class,
with the explicit compatible rational roots
$(-1)^{1/3^n}=-1$.  The local decomposition-section
comparison therefore changes no section class.

Let $X,Y$ be the logarithms of the standard tripod loops
and $C=[X,Y]$.  The $H$ lattice is
$\mathbb Z_3(3X)\oplus\mathbb Z_3(3Y)\oplus\mathbb Z_3 C$.
Wickelgren's evaluation of the standard path cocycle in
the subsection ``Evaluating the 2-nilpotent quotient of
$\mathfrak f$'', equation
$f(\sigma)=(\chi(\sigma)^2-1)/24$
\cite{BC_Wickelgren2012}, says that the section at $Q_c$
has central cocycle $fC$ relative to $P$.
The sign follows from the actual action convention
$g(y)=\mathfrak f_g^{-1}y^{\chi(g)}\mathfrak f_g$:
the lift $\mathfrak f_gs_P(g)$ conjugates $y$ to its
$\chi(g)$-th power.  Thus the $P,Q_c$ origins and their
peripheral generators are $0,fC$ and $3X,3Y$.

Consider the rational coordinate interchange
$(x,y,z)\mapsto(y,x,z)$ on the cubic.  It covers
$t\mapsto1-t$, interchanges $P,Q_c$, and fixes $O$.
Transporting the fundamental group along the same
standard path gives an actual automorphism of the
arithmetic $H$-extension.  Its geometric map swaps the
two abelian coordinates and negates $C$, and its change
of arithmetic origin is $fC$.  There are no additional
linear central terms in its geometric map: because
$fC$ is central, the geometric map commutes with the
diagonal $G_{\mathbb Q}$-action; a term from weight one
to weight two would satisfy $(\chi-\chi^2)c=0$ and
therefore is zero in the torsion-free lattice.
Consequently the exact automorphism is
\[
 ((u,v,w),g)\longmapsto((v,u,-w+f(g)),g).                   \tag{B.9.4}
\]
Its formula also checks directly that it is a continuous
group homomorphism and that its square is the identity.
It preserves $H$ because the covering subgroup is the
kernel of the two abelian coordinates modulo three.

At $O$, use the local parameter $s=z/x$.
The regularized leading terms of $u_0,v_0$ are $1,-1$;
thus the inherited abelian class of the unit tangent is
zero.  Conjugating by an element of $H$ makes its
abelian cocycle identically zero.  Its remaining class
is some $h\in H^1(G_{\mathbb Q},\mathbb Z_3(2))$.
The coordinate interchange takes $s$ to $z/y$, whose
leading coefficient at $O$ is $-1$.  The rational root
comparison above makes these two pro-three tangent
sections conjugate.  Since (B.9.4) acts on central
classes by $h\mapsto-h+[f]$, we obtain
$2h=[f]$.  Multiplication by two is invertible on this
group, so $h=[f/2]$.  A central geometric conjugation
then chooses the representative $f/2$.

We also check the inertia generator in this gauge.
Its abelian coordinates are $(-1,-1)$, since both
$u_0$ and $v_0$ have a simple pole in $s$.
Write its central coordinate as $w_0$.
The section just chosen is central, so its cyclotomic
normalization of the inertia generator and the
diagonal action give $\chi(g)^2w_0=\chi(g)w_0$ for
every $g$.  The lattice is torsion-free and the
character is nontrivial, hence $w_0=0$.
This determines the actual inertia image and avoids
an unsupported choice of a conjugate of $(xy)^{-3}$.
Twisting each of the three sections by its Kummer
parameter now gives (B.9.2), using the displayed group
law of E.29.  For each row the two abelian entries are
proportional, so that its BCH cross term is zero.
Equation (B.9.3) is the arithmetic map with exactly
the calculated image of the intrinsic reference
section; centrality of that image makes its geometric
part $q_O$ equivariant for the stated diagonal action.

Finally $H^1(G_{\mathbb Q},\mathbb Z_3(2))=\mathbb Z/3$
with generator $(\chi_3^2-1)/3$ by E.29m.
The coefficients $1/8$ and $1/16$ reduce to $2$ and
$1$ modulo three.  The three abelian packet lines
meet only at zero, and their central classes there
are distinct.  This proves all injectivity and
disjointness assertions.  The exact rational-arithmetic
program \texttt{code/bc\_h\_packets.py} verifies the
BCH automorphism, the three cocycle identities and
the three distinct central labels; its output is
included alongside the code.
\end{proof}

\begin{statement}{B.10}{Proved}{The full global image at the specified prime seven}
Let $A=N/\gamma_3N$, and use the normalized pair coordinates
$(a,b)$ and the $A^c$-action of A.15 for sections supported
at $\{3,7,\infty\}$.  Such a class has a full $N$-lift
exactly in the following cases:
\[
 b=1,
 \quad\text{or its restriction to }G_{\mathbb Q_7}
       \text{ belongs to one of the nine ramified packets (B.8.1).}
                                                               \tag{B.10.1}
\]
Equivalently these are precisely the classes which lift to
every lower-central quotient.  A full lift can be chosen
unramified outside $\{3,7,\infty\}$.

Here is an exact packet parametrization in the pair
coordinates.  Fix the embedding $\zeta_3\equiv2\pmod7$,
a Frobenius $\phi$ fixing compatible roots of $7$, and the
actual unit-tangent values $g_d=s_{d,1}(\phi)\in A$.
Use the same based inertia element $c_d\in A$ in each case.
The ramified part of the global image is
\[
 \left\{
  \bigl[(c_d^{2j}g_d,c_d^v)\cdot h\bigr]:
    d\in D(\mathbb Q_7),\quad
    v\in\mathbb Z_3\setminus\{0\},\quad j\in\{0,1,2\},
    \quad h\in A/A^c
 \right\}.                                                   \tag{B.10.2}
\]
Here $(a,b)\cdot h=(h^{-1}aF(h),h^{-1}bh)$, so $A/A^c$
means right cosets, and the displayed parameters are
injective after these explicit choices.  The symbols
$g_d$ retain the actual tangential Frobenius origins;
they are not replaced by arbitrary solutions of a
graded normalization equation.
\end{statement}
\begin{proof}
The exact isomorphism of arithmetic groups in A.14 and
the normalized-cocycle construction in A.15 identify
every global class with the asserted pair.
If $b=1$, lift $a$ arbitrarily to $\widetilde a\in N$
and take $\widetilde b=1$.  The full tame equation holds,
and A.15 extends this pair to the required global section.
If $b\ne1$ and a full section exists, B.7 forces its
local restriction to be cuspidal.  Conversely a local
cuspidal class in (B.8.1) has its actual full local lift.
Conjugate this lift by a lift in $N$ of the geometric
element comparing its class-two image to the specified
pair $(a,b)$.  Its generator values then project to
that exact pair.  A.15 extends them, with their conjugate
values on the second local factor, to a full normalized
global section.  This proves both directions of
(B.10.1), including the ramification support.
The nilpotent assertion follows either from A.15's
compactness proof or from B.11 below.

For (B.10.2), the completed local Kummer group is
$\mathbb Z_3\kappa(7)\oplus(\mathbb Z/3)\kappa(\zeta_3)$.
Choose compatible roots of $\zeta_3$.  Its unramified
cocycle is $(\chi_3-1)/3$, which has value two at
$\phi$ and value zero on tame inertia.
The uniformizer cocycle has values zero and one there.
The full decomposition-group formula of B.8 therefore
gives exactly the pair $(c_d^{2j}g_d,c_d^v)$.
Taking $j+3$ instead changes the pair by conjugation
with $c_d$, so the three chosen representatives cover
precisely the unit parameter group.

Different $d,v,j$ have different local section classes
by B.8.  Within each local class the geometric
centralizer in $A$ is trivial: its abelian image would
be fixed by Frobenius on $V$, and its remaining central
image would be fixed on $\bigwedge^2V$; the respective
eigenvalue absolute values are $\sqrt7,7$ and
$7,7^{3/2},49$, all greater than one.
Thus the full $A$-conjugacy orbit of that local pair
is freely parametrized by $h\in A$.
The normalized global conjugacies are exactly $A^c$
by A.15, acting on the right.  This proves the stated
right-coset parametrization and its injectivity.
\end{proof}

\begin{statement}{B.11}{Proved}{No new ramification above class two and compactness for a finite set of primes}
Let $S$ be any finite set of rational places containing
$3$ and infinity, and let $s_2$ be a class-two section
unramified outside $S$.  Every lift of $s_2$ to
$N/\gamma_rN$, $r\ge3$, is unramified outside the same
set $S$.  Moreover the following assertions are equivalent:
\[
 \begin{gathered}
 s_2\text{ lifts to }N;\\
 s_2\text{ lifts to }N/\gamma_rN\text{ for every }r\ge3.
 \end{gathered}                                               \tag{B.11.1}
\]
The separately given nilpotent lifts need not be
compatible.  The conclusion supplies a compatible
continuous full lift, also unramified outside $S$.
\end{statement}
\begin{proof}
For each $r$, iterated commutators give a continuous
$G_{\mathbb Q}$-equivariant surjection
\[
             V^{\otimes r}\longrightarrow
                    \gamma_rN/\gamma_{r+1}N.                 \tag{B.11.2}
\]
The target is a finite free $\mathbb Z_3$-module.
This integral assertion follows by Hall collection
in the free pro-three group: basic commutators of
weight $r$ span the quotient, and their Magnus
expansions have the linearly independent leading
Hall Lie monomials.  Collection at finite exponent
quotients, followed by compactness, gives the same
spanning assertion over $\mathbb Z_3$.
The leading monomials prove independence and exclude
torsion.  Thus (B.11.2) is the actual integral
graded comparison being used.

For $q\ne3$, Frobenius eigenvalues on $V$ have weights
one and two: on the elliptic factor their absolute
values are $\sqrt q$ and on the boundary factor they
are $q$.  The eigenvalues on (B.11.2) are therefore
among products of $r$ such eigenvalues.  For $r\ge3$
their absolute values are at least $q^{r/2}>q$.
This eigenvalue conclusion is valid for a quotient
of the tensor representation and does not assume
semisimplicity of the extension defining $V$.

Suppose a lift from class $r-1$ to class $r$ acquires
inertia at a prime outside $S$.  Its geometric
inertia values lie in the central kernel
$\gamma_rN/\gamma_{r+1}N$ and, as in B.7, factor
through the tame three-part.  The exact tame relation
requires its value $t$ to satisfy $F_qt=qt$;
conjugation by the section's Frobenius value acts
trivially on this central kernel.  For $r\ge3$ the
preceding eigenvalue calculation and torsion-freeness
give $t=0$.  Induction starting at the fixed class-two
section proves that every higher nilpotent lift has
the asserted support.

Let $K_0=\mathbb Q(\zeta_3)$ and let $S'$ be its
places over $S$.  The relative pro-three quotient
$\Gamma_S$ of $G_{\mathbb Q,S}$ has kernel
$G_{K_0,S'}(3)$ and quotient of order two.
That pro-three kernel is finitely generated: its
$\mathbb F_3$-characters form the finite Kummer group
of $S'$-units and the finite three-primary ideal
class group.  The finite-dimensional Frattini
quotient generates a pro-three group, by the same
finite-quotient maximal-subgroup argument used in A.14.
Thus $\Gamma_S$ is finitely generated.
The actual geometric action factors through it, as
proved in A.15 from the three-group action on the
Frattini quotient of $N$.  Every cocycle under
consideration also factors through $\Gamma_S$:
over $K_0$ its graph is a homomorphism to a
pro-three semidirect product.

Fix generators $g_1,\ldots,g_m$ of $\Gamma_S$ and
a cocycle representative of $s_2$ on it.
Consider the compact space of tuples in $N^m$
whose projections are its fixed generator values.
Choose a free profinite presentation on these
generators, with its closed subgroup of relators.
For each $r\ge3$, impose all those relators on
the lifted tuples in
$(N/\gamma_rN)\rtimes\Gamma_S$.
These are closed conditions: at any finite quotient
each profinite word is a continuous function of
the generator values.  A lift at level $r$ makes
the resulting closed subset nonempty, after
conjugating its projection to the fixed representative;
the required geometric conjugator lifts because
$N/\gamma_rN\to N/\gamma_3N$ is surjective.
The closed sets are nested.  Compactness supplies
one tuple satisfying every level of every relator.

The intersection $\bigcap_r\gamma_rN$ is trivial:
a nontrivial element survives in a finite three-group,
and that group is nilpotent.  Thus the tuple satisfies
all relators in $N\rtimes\Gamma_S$ itself and defines
a continuous section.  Composing with
$G_{\mathbb Q,S}\to\Gamma_S$ gives the desired full
cocycle with the prescribed class-two projection.
This proves the nontrivial implication in (B.11.1);
the other is induced by the continuous quotient maps.
\end{proof}

\begin{statement}{B.2}{Refuted}{The inherited class-two image need not be a point or cusp image}
The exclusion proposed in Lead 1 is false.  There are sections of
$N\rtimes G_{\mathbb Q}$ whose image in $\operatorname{Sec}(H)$
has toric coordinates $(\kappa(3),2\kappa(3))$.  Every such image
lies outside the three rational cusp packets, and $U(\mathbb Q)$
has no point producing it.  It therefore supplies the asserted
counterexample already with a full geometric pro-$3$ lift.
\end{statement}
\begin{proof}
We use exactly the onto map in LM7 E.34: the coordinate-unit
homomorphism on the unramified group
$H^1(G_{\mathbb Q,\{3,\infty\}},V)$ has image
$\mathbb Z_3\kappa(3)^2$, and every class in this source lifts to
the full $N$.  Its projectivity step is checked as follows.
Over $\mathbb Q(\zeta)$, the maximal pro-$3$ Galois group unramified
away from three is free on two generators: Kummer theory gives
the two classes $\zeta,1-\zeta$; the class group of
$\mathbb Z[\zeta,1/3]$ is zero, and the three-primary Brauer group
supported at its single omitted finite place is zero by the sum of
local invariants.  Thus degree-two cohomology with $\mathbb F_3$
vanishes.  The free-pro-$3$ criterion follows by solving finite
central embedding problems.  Complex conjugation has order two,
so averaging gives the same degree-two vanishing for the relative
semidirect product.  The geometric action and every abelian cocycle
factor through this relative group: their restrictions over
$\mathbb Q(\zeta)$ take values in pro-$3$ groups.  Successive finite
geometric embedding problems therefore have solutions; finite
generation of the relative group and compactness make them
compatible.  This is the full-group lifting step used from E.34.
The direct-summand coordinate-unit map is the explicit algebraic
section $(a,b)\mapsto(a,b,1,1,1)$ of (E24.4), so its onto property
requires no further lifting assertion.  These arguments also agree
with the free-Galois-group mechanism in Hoshi, Theorem 3.5
\cite{BC_Hoshi2010}; that theorem concerns the specified proper
Jacobian sections, while the preceding relative-group argument
is the needed one for all ten coordinates here.

By (B.4), the cusp projections are $(c,0)$, $(0,c)$ and
$(-c,-c)$, with $c\in H^1(G_{\mathbb Q},\mathbb Z_3(1))$.
The valuation at three sends the proposed vector to $(1,2)$,
which lies on none of the three lines, already over $\mathbb F_3$.
The code verifies this finite separation.  Conjugation in $H$
does not alter the abelian cohomology class, so no choice of
central lift can change this conclusion.  LM7's computation
$E(\mathbb Q)=\{O,P,Q_c\}$ gives $U(\mathbb Q)=\varnothing$.

The drafter's suggested vector
$\kappa(3)(e_O+e_P)$ does have an intrinsic abelian image outside
the packets (B.1): its elliptic image is zero and its toric image
is not diagonal.  Its inherited toric image is
$(0,-\kappa(3))$, on the $Q_c$ line.  That particular toric
calculation therefore gives no exclusion in $\operatorname{Sec}(H)$.
The vector used above supplies the required exclusion there.
\end{proof}

\begin{statement}{B.3}{Proved}{The cusp packets in the cuspidally central quotient}
In $N^{\rm cc}=C\times_TV$, let the generator of the central
$\mathbb Z_3(1)$ of $C$ be the elliptic commutator $[a,b]$.
Choose the reference $C$-section above the elliptic class zero
to be the tangent at $O$; above $\kappa(P-O)$ and
$\kappa(Q_c-O)$ choose the images of the respective unit tangents.
Use these references for every abelian class in the corresponding
three elliptic fibers.  With these explicit choices for the
torsors denoted by $s(\alpha,a)$, the cusp packets are
\[
 \begin{array}{c|c|c}
 d&\alpha&a\\ \hline
 O&e_O\kappa(u)&-\kappa(u)\\
 P&\alpha_P+e_P\kappa(u)&0\\
 Q_c&\alpha_Q+e_{Q_c}\kappa(u)&0
 \end{array}                                                \tag{B.6}
\]
Here $u$ ranges over the full pro-$3$ completion of
$\mathbb Q^\times$.  This parametrization is injective on each
row, and the rows have disjoint images.
\end{statement}
\begin{proof}
The morphism $U\hookrightarrow E\setminus\{O\}$ fills the cusps
$P,Q_c$, so their inertia maps to the identity in $C$.  At $O$,
the once-punctured elliptic presentation gives
$c_O=-[a,b]$ in the central lattice of $C$.
Consequently a cusp twist $\kappa(u)$ changes the $C$-coordinate
by $-\kappa(u)$ at $O$ and by zero at $P,Q_c$.
The $V$-coordinate is exactly (B.1).
The map to $V$ and its already proved packet disjointness and
parameter injectivity give both remaining assertions.
\end{proof}

\begin{statement}{B.4}{Proved}{Lead 4 for the full geometric pro-three group}
For every continuous section
$s:G_{\mathbb Q}\to N\rtimes G_{\mathbb Q}$ put
$\Delta=\ker(\chi:G_{\mathbb Q}\to\widehat{\mathbb Z}^{\times})$.
Then
\[
 C_N(s(\Delta))=
 \begin{cases}
 I_{\widetilde d},&s\text{ is cuspidal at the lifted cusp }\widetilde d,\\
 1,&s\text{ is noncuspidal}.
 \end{cases}                                                \tag{B.7}
\]
In the first case the lifted cusp is unique and its image is one
of $O,P,Q_c$.  Thus both implications in Lead 4 hold for the
actual pro-$3$ group, with its genuine section action.
\end{statement}
\begin{proof}
We give the pro-$3$ argument.  No fixed element is lifted from a
quotient of a full profinite fixed group.

For an open subgroup $H\subset N$, the corresponding finite
geometric cover has maximal pro-$3$ fundamental group $H$.
To check this comparison, choose an open normal subgroup
$M\subset H$ of $N$; $N/M$ is a finite $3$-group and $H/M$ is
subnormal in it.  Pulling back a subnormal chain expresses the
cover as a tower of covers of $3$-power degree with normal
successive groups.  A finite $3$-group quotient of a subgroup
at one stage becomes a quotient of a finite $3$-group after
taking the finitely many conjugates and then the extension by
the stage's $3$-group.  Induction proves that the topology
induced from the original maximal pro-$3$ quotient is the full
pro-$3$ topology on the subgroup.  This proves the comparison.
For its compactification $X_H$ and geometric boundary $D_H$,
the oriented surface presentation gives the exact sequence
\[
 \mathbb Z_3(1)\xrightarrow{\delta}
 \mathbb Z_3(1)[D_H]\longrightarrow H^{\rm ab}
 \longrightarrow T_3J_{X_H}\longrightarrow0.                \tag{B.8}
\]
Every arrow is a continuous homomorphism, and the last two
arrows are induced respectively by inertia and by filling the
cusps.  Its only boundary relation is the diagonal sum.

Each lifted inertia group is $\mathbb Z_3(1)$.  Two distinct
lifted cusps have inertia groups meeting trivially.  Indeed,
separate the cusps in a finite normal pro-$3$ cover, refine
until there are at least three boundary points, and take a
positive $3$-power of a purported common element so that it
lies in the cover group.  It remains nonzero in
$\mathbb Z_3(1)$ and (B.8) puts it on two different boundary
lines; those lines have zero intersection.  The refinement
exists because an inertia group has infinite index: an open
procyclic subgroup would give a hyperbolic cover whose
abelianization has rank one, contradicting its surface
presentation.  Normal covers therefore have unboundedly
many points above a cusp.  The same intersection argument
shows that the geometric normalizer of nontrivial inertia is
that inertia group, and its arithmetic normalizer is the
corresponding cusp decomposition group: a normalizer fixes
the lifted cusp, and stabilizers are precisely those groups.
These geometric assertions also agree with Nakamura's
pro-class statements, Lemmas (2.1.2) and (2.1.4)
\cite{BC_Nakamura1994}.

Let $z\in C_N(s(\Delta))$.  The action through $s$ on the
finitely generated pro-$3$ group $N$ has an open stabilizer
for each open subgroup $H$: a characteristic open subgroup
inside $H$ reduces the assertion to an action on a finite
group.  Consequently $H$ is realized by a cover over some
number field $L$, and $s|_{G_L}$ lifts to that cover when
$H$ contains $z$.  The image of $z$ in (B.8) is fixed by
$G_{L(\mu_\infty)}$.
Ribet's Appendix Theorem 1 \cite{BC_Ribet1981} says that
an abelian variety over a number field has finite torsion
over its full cyclotomic extension.  Its consequence here
is exact: a nonzero invariant vector of $T_3J_{X_H}$ would
give cyclotomic rational torsion points of unbounded
$3$-power order.  Thus the image of $z$ in the last term
of (B.8) is zero.

Apply this to $H=M\langle z\rangle$, for every open normal
$M\subset N$.  The image of $z$ in $H^{\rm ab}$ belongs
to the boundary subgroup.  The quotient $H/M$ is cyclic of
$3$-power order and generated by the image of $z$.
Its boundary images generate it, and a cyclic $3$-group
generated by subgroups equals one of those subgroups.
Thus $z$ fixes a cusp in the cover associated to $M$.
The finite sets of such fixed cusps are nonempty at every
normal level.  A common refinement supplies every finite
set of compatibility requirements; compactness gives a
lifted cusp $\widetilde d$ fixed by $z$.  Hence
$z\in I_{\widetilde d}$.

We verify that all nonzero elements of the centralizer
belong to one inertia group.  More generally let a closed
subgroup $B\subset N$ consist entirely of inertia elements.
The set of cusps with $B\cap I_c\ne1$ is $B$-stable.  If
it were a single orbit while every $B\cap I_c$ were proper,
$B$ would be a union of conjugates of one proper closed
subgroup.  A finite quotient in which this subgroup remains
proper contradicts this: in the finite coset action the
average number of fixed cosets is one and the identity
fixes more than one, so some element fixes none.
We can therefore choose two distinct $B$-orbits of cusps
unless $B$ is contained in one inertia group.  Compactness
separates these orbits at a finite normal level $M$; the
open subgroup $BM$ has two distinct corresponding cusps.
When it has only two cusps, its compactification has
positive genus.  Every element of $B$ maps trivially to
the compactification group.  A degree-three finite
unramified cover of that compactification consequently
gives a further open subgroup containing $B$ with six
cusps.  We have thus arranged at least three cusps while
keeping the specified two distinct.

Take nonzero $z,z'$ in the two inertia intersections.
Their boundary images are $ae_1,be_2$, with $a,b\ne0$.
Both $zz'$ and $z^2z'$ have nonzero boundary image.
With four or more cusps, neither image can be on a
single boundary line.  With three cusps, lying on one
line forces respectively $a=b$ and $2a=b$, a contradiction.
This contradicts the defining property of $B$ and proves
the desired single-inertia assertion.

For $C=C_N(s(\Delta))\ne1$, we have proved
$C\subset I_{\widetilde d}$.  Each $s(\sigma)$,
$\sigma\in\Delta$, fixes a nonzero element there;
the inertia-intersection result forces it to fix
$\widetilde d$.  Its inertia action is $\chi_3(\sigma)=1$,
so $I_{\widetilde d}\subset C$.  Thus equality holds.
Normality of $\Delta$ makes $C$ stable under all of
$s(G_{\mathbb Q})$.  The arithmetic normalizer identity
puts the entire section in the decomposition group of
$\widetilde d$.  Its base cusp is rational, and uniqueness
follows from inertia intersections.  Conversely a
cuspidal section fixes its inertia over $\Delta$;
the same argument gives equality in (B.7).
\end{proof}

\begin{statement}{B.5}{Proved}{The cuspidal members of the full unramified family}
Let $\mathscr F_S$ be the image of
$H^1(G_{\mathbb Q,\{3,\infty\}},N)$ in
$H^1(G_{\mathbb Q},N)$, using the integral tangent at $O$.
Its cuspidal subset is the disjoint union of the three
packets obtained by twisting the unit tangents at
$O,P,Q_c$ by $t\kappa(3)$, $t\in\mathbb Z_3$.
Every other full section class in $\mathscr F_S$ has
centralizer $1$ over $\Delta$.  The centralizer of each
cuspidal member is its complete pro-$3$ inertia group.
\end{statement}
\begin{proof}
The good marked model over $\mathbb Z[1/3]$ and its
integral nonzero tangents give actions and tangential
cocycles unramified outside three.  This is the based
smooth-specialization step used in E.34; the comparison
is made on covers with their local parameter root
fiber functors, so it applies to the genuine action.
For a twist $u\in\mathbb Q^{\times,\wedge_3}$,
restriction to inertia at $q\ne3$ has abelian value
$v_q(u)e_d$.  The vector $e_d$ is primitive and nonzero
in $V$, so a section in $\mathscr F_S$ forces
$v_q(u)=0$.  Unique factorization in the completed
Kummer group then gives $u=3^t$, with $t\in\mathbb Z_3$.
Conversely these twists are unramified outside three.
Statement B.4 identifies every cuspidal full section
with exactly one of these packets and computes all
remaining centralizers.  The abelianized images are
the three affine lines in B.1 with the parameter
restricted to $t\kappa(3)$.

This assertion concerns full section classes.
The surjection from this family to its abelianized
classes in E.34 need not have one-point fibers.
An abelian class outside the three displayed lines
has exclusively noncuspidal full lifts.  We have not
deduced cuspidality of every lift above a point on one
of those lines.
\end{proof}

\begin{statement}{B.12}{Proved}{The marked local classification when the filled-curve projection is unramified}
Let $k/\mathbb Q_q$ be finite, $q\ne3$, with residue cardinality
$Q$, and let $d=Q_c$ be the marked cusp. Write
\[
 \mathscr Q=N/\gamma_3I_{\rm mark},\qquad
 B=N/I_{\rm mark},\qquad
 \mathscr I=I_{\rm mark}/\gamma_3I_{\rm mark},\qquad
 Z=\gamma_2I_{\rm mark}/\gamma_3I_{\rm mark}.
 \tag{B.12.1}
\]
All groups carry the actual arithmetic action, and the
quotient $B$ is the full geometric pro-three group of the
curve obtained by filling $d$.
Choose tame three-inertia $\tau$, arithmetic Frobenius
$\phi$, the peripheral generator $c$ at $d$, and the
unit tangential reference $s_{d,1}$ compatibly.
Put $F_0=s_{d,1}(\phi)$.
The induced based action $\Phi$ of $F_0$ on $B$ is
the action of the point $d$ on the filled curve.
There are actual equivariant identifications
\[
 M:=\mathscr I^{\rm ab}
     =\mathbb Z_3[[B]](1),\qquad
 Z=\widehat{\bigwedge}^{\,2}M,\qquad
 (F_0)_M=Q\Phi,\quad (F_0)_Z=Q^2\widehat{\bigwedge}^{\,2}\Phi.
 \tag{B.12.2}
\]
Here the completed modules are inverse limits over finite
continuous partitions of the profinite label set $B$;
the label $1$ is the chosen lift of the marked cusp.

Every local section of $\mathscr Q$ whose projection to $B$
is unramified and whose tame inertia is nontrivial has,
up to geometric conjugation, exactly the following form:
\[
 s(\tau)=c^a,\qquad s(\phi)=c^r\exp(z)F_0,
 \quad
 \begin{cases}
 a\in\mathbb Z_3\setminus\{0\},\\
 r\in\mathbb Z_3/(Q-1)\mathbb Z_3,\\
 [z]\in\mathcal C_{k,d}:=
 Z/(Q^2\widehat{\bigwedge}^{\,2}\Phi-1)Z .
 \end{cases}                                                \tag{B.12.3}
\]
In particular its $B$-section is the point section at $d$.
The quotient in (B.12.3) uses the actual image; that image
is closed because $Z$ is compact.
Every displayed parameter occurs, and two displayed
sections are conjugate exactly when all three parameters
agree.

Under $\mathscr Q\to N/\gamma_3N$, their image is exactly
the ramified local cusp packet at $d$. For each member of
that packet the entire local fiber is the group
$\mathcal C_{k,d}$, acting simply transitively after the
displayed reference is fixed.
Finally a local marked section whose abelian tame inertia
in $V$ is zero is itself unramified.
\end{statement}
\begin{proof}
The filling map gives a geometric free product
$N=B*_3\langle c\rangle$ after a choice of free lift
of $B$. Its normal kernel is the free pro-three group on
the profinite set of lifted marked points, a torsor under
$B$. The two inverse semidirect-product maps used in
F.25, now with $b\mapsto bcb^{-1}$, prove this assertion
for the actual filling kernel. The maximal class-two
quotient consequently has abelianization
$\mathbb Z_3[[B]](1)$ and center its completed exterior
square. Finite partitions and the finite-rank Hall basis
prove these assertions and their inverse-limit
compatibility. Galois permutes the labeled inertia
generators, up to conjugacy and cyclotomic power.
This proves (B.12.2); at the distinguished label the
unit tangent also gives the exact equality
$F_0cF_0^{-1}=c^Q$.

We record the fixed-point and measure facts needed below.
For any unramified $B$-section, every positive power of
its Frobenius automorphism has fixed subgroup one in $B$.
On $B^{\rm ab}$ the eigenvalues have complex absolute
values $Q$ or $\sqrt Q$; the same assertion for the
degree-$r$ lower-central quotient follows from its
equivariant Hall injection into $(B^{\rm ab})^{\otimes r}$.
Inner twists do not change the action on this graded
quotient. Thus no such power has eigenvalue one in any
positive degree. The intersection of the closed lower
central groups is one, as finite three-group quotients
are nilpotent. This proves the fixed-subgroup assertion
for the full $B$.
It follows that an affine Frobenius action on a $B$-torsor
has at most one point fixed by any given positive power.
A finite orbit must therefore be a single fixed point,
and there is at most one such point in total.

The closure of that affine action is virtually pro-three.
Indeed the translation subgroup is pro-three and the
kernel of the action on the Frattini quotient of the
finitely generated free pro-three group is pro-three.
An integer power removes the finite prime-to-three part.
For an invariant $\mathbb Z_3$-measure on such a torsor,
all mass outside its possible single periodic point is
zero. The same assertion holds with $\mathbb F_3$
coefficients. To see this precisely, stabilize a clopen
set avoiding the periodic point by an open subgroup
of the pro-three closure. Every orbit there is infinite.
For any $n$, compactness supplies a finite invariant
partition on which every orbit has size divisible by
$3^n$. Pushforward sums make the measure of the
clopen set divisible by $3^n$; over $\mathbb F_3$
one step already makes it zero.
This is the integral measure argument of F.25, applied
to the actual permutation action.

Now let the projection of $s$ to $B$ be unramified.
Wild inertia and the tame factors other than three
map trivially to the geometric pro-three group.
Write $t=s(\tau)\in\mathscr I$ and $F=s(\phi)$,
so $FtF^{-1}=t^Q$.
On $M$ the action of $F$ is $Q\Psi$ for the affine
Frobenius on the marked-point path torsor. The image
$m$ of $t$ therefore satisfies $\Psi m=m$.
The measure argument makes $m=a\delta_x$ for an actual
fixed label $x$, or makes it zero when there is no such
label.

When $m=0$, the element $t$ is in $Z$.
The fixed relation says
$(Q\widehat{\bigwedge}^{\,2}\Psi-1)t=0$ after dividing
by the unit $Q$. This operator is injective.
On every finite label quotient the permutation has
finite order, so all its eigenvalues are roots of
unity; none becomes one after multiplication by
the rational integer $Q>1$. Tensoring with $\mathbb Q_3$
therefore proves injectivity on that finite free lattice.
The inverse-limit comparison in (B.12.2) then proves
injectivity on $Z$. Hence $t=1$.

For nontrivial $t$ we consequently have $a\ne0$.
The fixed path $x$ is exactly an isomorphism between
the projected section and the point section at the
filled point $d$: on inertia both actions are trivial
by good marked reduction, and on the unramified
quotient it is the displayed Frobenius fixed-path
condition. Conjugate by a lift of that path in
$\mathscr Q$. The label becomes $1$, the $B$-projection
is the specified point section, and $F=hF_0$ with
$h\in\mathscr I$.

Use the integral class-two logarithm, with its sole
denominator two, and set $e=\delta_1$.
Write
\[
                   t=(ae,z),\qquad b=h_{\rm ab}.
 \]
The equality $F_0cF_0^{-1}=c^Q$ removes any central
origin term on the line $\mathbb Z_3e$.
Class-two collection now gives the exact equation
\[
       (Q\widehat{\bigwedge}^{\,2}\Phi-1)z
                           =a\,e\wedge b.                  \tag{B.12.4}
\]
In fact $F_0t=(Qae,Q^2(\bigwedge^2\Phi)z)$ and
conjugating by $h$ adds $aQ\,b\wedge e$; comparison
with $t^Q=(Qae,Qz)$ gives (B.12.4).

Augmentation and the fixed label give the equivariant
splitting $M=\mathbb Z_3e\oplus M^0$.
Projection of (B.12.4) to
$\widehat{\bigwedge}^{\,2}M^0$ is zero, by the same
finite-permutation eigenvalue argument. Thus
$z=e\wedge v$ for a unique $v\in M^0$, and
\[
                     (Q\Phi-1)v=a b^0,
           \qquad b^0=b-\operatorname{aug}(b)e.             \tag{B.12.5}
\]
The operator $Q\Phi-1$ on $M^0$ is injective modulo
three. If $Q=1\bmod3$, its kernel consists of
$\Phi$-invariant $\mathbb F_3$-measures; if
$Q=-1\bmod3$, every element of its kernel is fixed
by $\Phi^2$. Either way the measure argument supports
it only at $1$, and augmentation zero removes that
last possibility. Reduction of the completed module
commutes with the finite label inverse limits:
compatible lifts exist by compactness, or by the
surjective transition maps of the finite free
modules.
Writing $a=3^s u$ with $u$ a unit, repeated reduction
of (B.12.5) now gives $v\in3^sM^0=aM^0$.
For a lift $h_1$ of $v/a$, the explicit conjugation
$t\mapsto h_1t h_1^{-1}$ changes $t$ to $c^a$.
This is an integral divisibility argument; no
division by a nonunit has been assumed.

An element of $\mathscr Q$ centralizing $c^a$ must
have trivial $B$-image, because the action of $B$
on the label torsor is free and
$a\delta_b=a\delta_1$ forces $b=1$.
Inside $\mathscr I$, the commutator equation is
$a e\wedge h_{\rm ab}=0$. Torsion-freeness and the
augmentation splitting give $h_{\rm ab}\in\mathbb Z_3e$.
Consequently the exact centralizer is
\[
                         C_{\mathscr Q}(c^a)
                                  =\langle c\rangle Z.     \tag{B.12.6}
\]
The relation with $F_0$ puts $F F_0^{-1}$ in this
centralizer, giving precisely $c^r\exp(z)F_0$.
Conjugating by $c^u\exp(z')$ changes the two
Frobenius parameters by
$(Q-1)u$ and $(Q^2\bigwedge^2\Phi-1)z'$.
There is no cross term, since $Z$ is central in
$\mathscr I$ and $c$ commutes with it.
Abelianization fixes $a$, and (B.12.6) proves that
these are all possible identifications.
Conversely every displayed pair satisfies the
tame relation. Continuous powering gives the
homomorphism on the procyclic unramified quotient,
and the tame semidirect-product presentation gives
an actual continuous local section. This proves
(B.12.3).

All conjugates of the marked generator have the
same image $e_d$ in $V$. Their pairwise commutators
therefore vanish in $\gamma_2N/\gamma_3N$.
Thus the actual quotient homomorphism sends
$Z=\gamma_2I_{\rm mark}/\gamma_3I_{\rm mark}$ to one
in $N/\gamma_3N$. Formula (B.12.3) has exactly the
stated class-two cusp image.
Conversely, any marked section above a ramified
class-two cusp at $d$ has zero abelian inertia in
the filled group. The full local theorem B.7 for
that filled curve makes its $B$-projection unramified,
so the preceding classification includes its entire
fiber.
The same argument starts from any marked section
with zero inertia in $V$: its $B$-projection is
unramified, while $a e_d=0$ forces $a=0$.
The already proved zero-$m$ case then makes $t=1$.
\end{proof}

\begin{statement}{B.13}{Proved}{The central marked local fiber is a countable product of three-adic integers}
The group $\mathcal C_{k,d}$ in (B.12.3) is a nonzero,
torsion-free, compact metrizable $\mathbb Z_3$-module.
After a choice of divisible basis of its discrete dual,
there is a continuous $\mathbb Z_3$-linear isomorphism
\[
                         \mathcal C_{k,d}
                                  \simeq\prod_{n\ge1}\mathbb Z_3 .
                                                               \tag{B.13.1}
\]
The coordinates in (B.13.1) depend on that stated choice;
the quotient and the section action in (B.12.3) are
intrinsic to the displayed reference.
In particular there are actual local marked sections
with the same ramified class-two cusp image which
are not cuspidal and have no full-$N$ lift.
\end{statement}
\begin{proof}
Let $D=Q^2\bigwedge^2\Phi-1$.
The completed exterior square identifies with the
antisymmetric direct summand of
$\mathbb Z_3[[B\times B]]$ by the idempotent
$(1-\mathrm{swap})/2$.
Modulo three, $D$ is $\bigwedge^2\Phi-1$.
The only periodic point of $\Phi$ on $B$ is $1$,
so the only periodic point on $B\times B$ is
$(1,1)$. The invariant-measure argument of B.12
makes every invariant $\mathbb F_3$-measure a
multiple of $\delta_{(1,1)}$. Its antisymmetric
part is zero. Thus $D$ is injective modulo three.

If $3z=Dv$, reduction modulo three gives $v=3v_1$,
and cancellation in the torsion-free module $Z$
gives $z=Dv_1$. This proves that the cokernel has
no three-torsion, and hence no torsion at all.
It is compact and metrizable because $B$ is a
finitely generated pro-three group, its finite
characteristic quotients form a countable cofinal
family, and all the coefficient quotients are
finite. The image $DZ$ is compact and closed.

We prove infinite rank with actual finite quotient
maps. Put $V_B=B^{\rm ab}$, of rank nine, and
$\Omega_n=V_B/3^nV_B$.
The homomorphism $B\to\Omega_n$ induces a continuous
surjection
\[
 Z/3Z\longrightarrow\bigwedge^2\mathbb F_3[\Omega_n].
 \]
It sends $D$ to $\bigwedge^2\Phi_n-1$, and therefore
induces a surjection from $\mathcal C_{k,d}/3$ onto
the coinvariants of that finite permutation module.
If $m$ is the order of $\Phi$ on $V_B/3$, its order
on $\Omega_n$ divides $m3^{n-1}$: expand the powers
of the integral matrix $\Phi^m=1+3A$.
The exterior basis vectors
$\delta_0\wedge\delta_v$, $v\ne0$, form an invariant
direct summand. No orbit reverses its orientation,
because $\Phi_n$ fixes $0$ and does not send a
nonzero vector to $0$.
Their coinvariants have one independent vector
for each orbit, and consequently dimension at least
\[
                      \frac{3^{9n}-1}{m3^{n-1}} .
 \]
This tends to infinity. A concrete nonzero test
already at $n=1$ takes a rational unmarked cusp
vector $v$ and sums the coefficients on
$\delta_0\wedge\delta_v$ and, when distinct, on
$\delta_0\wedge\delta_{-v}$.
Frobenius multiplies that cusp vector by $Q$,
so the sum kills $D$ modulo three and takes the
value one on an actual central commutator
$[c,bcb^{-1}]$ whose label $b$ maps to $v$.

For an explicit description of the isomorphism in
(B.13.1), let $\mathcal C^\vee$ be the group of
continuous characters of $\mathcal C_{k,d}$ into
$\mathbb Q_3/\mathbb Z_3$.
It is countable and three-primary.
Injectivity of multiplication by three on
$\mathcal C_{k,d}$ makes multiplication by three
surjective on $\mathcal C^\vee$, so the latter
is divisible. This dual assertion can be checked
without a continuity assumption left implicit:
a finite-valued character on the closed subgroup
$3\mathcal C_{k,d}$ factors through its image in
a finite quotient of $\mathcal C_{k,d}$, and extends
there because $\mathbb Q_3/\mathbb Z_3$ is divisible.
Pulling that finite character back supplies the
continuous extension.
Choose a basis of its subgroup killed by three,
and compatible successive division chains for
these basis elements. Independence and exhaustion
follow by induction on the order of a character;
this gives an isomorphism
$\mathcal C^\vee\simeq\bigoplus_{n\ge1}
 \mathbb Q_3/\mathbb Z_3$.
The index set is countably infinite by the
unbounded finite dimensions just proved.
Evaluation on those chosen division chains is
the continuous homomorphism in (B.13.1).
Duality, or separation of points by finite
characters followed by compactness, proves
that it is an isomorphism.

The cuspidal class with the parameters $a,r$
has $[z]=0$. By B.12 a nonzero $[z]$ cannot
be removed by any geometric conjugation, so
it is not cuspidal. A full local lift would be
cuspidal by B.7, and would retain that property
under the actual quotient to $\mathscr Q$.
Thus these sections have no full lift.
\end{proof}

\begin{statement}{B.14}{Proved}{The complete local marked image and the unmarked ramified fibers}
Retain the local data and actual quotients of B.12.
If the projection of a section of $\mathscr Q$ to $B$
is ramified, that section is itself cuspidal at a
unique unmarked cusp $d'\in D(k)\setminus\{d\}$.
For a fixed ramified $B$-section the lift to
$\mathscr Q$, up to geometric conjugation, is unique
and is the image of its actual full cusp section.

Consequently the complete local image of
$\operatorname{Sec}_k(\mathscr Q)$ in
$\operatorname{Sec}_k(N/\gamma_3N)$ equals the
local full image: all unramified classes and the
disjoint ramified cusp packets of B.8.
The fiber over a ramified unmarked cusp class is
a singleton; the fiber over a ramified marked
cusp class is exactly the group $\mathcal C_{k,d}$
of B.12--B.13.
\end{statement}
\begin{proof}
By the full local theorem B.7 applied to the
filled curve, a ramified projected section is
cuspidal at an unmarked rational cusp $d'$.
Filling $d$ is an isomorphism on the inertia
and tangent data at $d'$. Thus this projected
section has an actual full cusp lift, which we
choose as reference. It remains to prove that
every lift to $\mathscr Q$ is conjugate to that
reference.

Let $X$ be the profinite torsor of paths on the
filled curve from this reference at $d'$ to the
filled point $d$. In these coordinates
$M=\mathbb Z_3[[X]](1)$.
The geometric inertia image of the reference
is $J_0=\langle c_{d'}^a\rangle\simeq\mathbb Z_3$,
with $a\ne0$. It acts freely on $X$ by translation,
and diagonally on $X\times X$.
Both are pro-three actions without finite orbits.
The integral measure lemma therefore gives
$M^{J_0}=0$ and
$(\mathbb Z_3[[X\times X]])^{J_0}=0$.
Taking the antisymmetric direct summand gives
$Z^{J_0}=0$ as well.

We justify continuous cohomology and orbit
coinvariants with these compact coefficients.
For a procyclic three-group generated by $T$
acting continuously on a compact pro-three
module $A$, evaluation at $T$ gives
\[
              H^1_{\rm cts}(J_0,A)=A/(T-1)A.               \tag{B.14.1}
\]
Every value $x\in A$ defines a continuous cocycle:
its norm sum at $T^{3^n}$ tends to zero in every
finite coefficient quotient, since on that
quotient the action has three-power order
and the repetitions are eventually divisible
by its exponent. These finite cocycles are
compatible and give the stated continuous
cocycle. Coboundaries are exactly $(T-1)A$;
this image is compact and closed.
In the permutation module case, pushforward
of measures gives an actual isomorphism
\[
 M/(T-1)M
               \simeq\mathbb Z_3[[J_0\backslash X]](1).
                                                               \tag{B.14.2}
\]
At each finite continuous invariant partition,
permutation coinvariants are the free module
on its orbits. Inverse limits preserve this
description: the finite systems of correcting
elements have compact solution sets, so the
kernel of pushforward is the actual closed
image $(T-1)M$. Equivalently the universal
continuous homomorphism from the completed
free module to an orbit-constant module has
exactly this kernel.
The same construction for $X\times X$ commutes
with the idempotent $(1-\mathrm{swap})/2$,
and hence computes $Z/(T-1)Z$ as the corresponding
antisymmetric direct summand of the orbit module.

Let $\Psi$ denote the actual Frobenius permutation
on $J_0\backslash X$.
The induced Frobenius action on (B.14.1) contains
an inverse factor $Q$: conjugation sends
$T$ to $T^Q$, so the pullback of a cocycle
evaluated at the generator is multiplied by
$Q^{-1}$ in coinvariants.
On $M$ the coefficient action has factor $Q$.
Thus the action on (B.14.2) is precisely $\Psi$,
with no remaining Tate factor.

There are no finite $\Psi$-orbits.
We construct the map proving this assertion.
Filling all remaining punctures maps the full
pro-three group of the filled curve onto $T_3E$.
On path torsors it gives a continuous equivariant
map from $X$ to the proper path torsor from
$d'$ to $d$. The group $J_0$ maps to one in
$T_3E$, because it is peripheral inertia.
The map consequently descends to
$J_0\backslash X$.
A point there periodic under $\Psi^n$ would
give a fixed point of the proper path torsor
under Frobenius to the $n$th power.
Original arithmetic inertia already acts
trivially on this torsor by good reduction.
Such a fixed point would make the completed
Kummer class of $d-d'$ zero over the finite
unramified degree-$n$ extension of $k$.

That difference is a nonzero three-torsion
point of $E$. At every finite extension of
a local field of residue characteristic
different from three, the reduction kernel
is pro-residue-characteristic and the
three-primary completion of $E(k')$ is its
finite three-primary torsion subgroup.
In particular its nonzero three-torsion
point remains nonzero in that completion.
The inverse-limit Kummer injection then
makes its class nonzero in $H^1(k',T_3E)$.
Being divisible once, or several times,
in a larger finite group would not make
this completed class zero. This proves
the required absence of periodic points.

The Frobenius action on the compact quotient
set is virtually pro-three, as in B.12.
Its invariant integral measure group is
therefore zero. Equations (B.14.1)--(B.14.2)
show
\[
             H^1_{\rm cts}(G_k,M)=0.                       \tag{B.14.3}
\]
For precision, the unramified quotient causes
no additional term: its possible first
cohomology term has coefficient $M^{J_0}=0$.
One can obtain the same conclusion without
invoking a spectral sequence. If the inertia
cocycle is a coboundary, remove it. The tame
relation then makes the Frobenius value
$J_0$-invariant, hence zero. Wild and
non-three tame inertia have no nonzero
homomorphism to these compact pro-three
modules.

For $Z$ the inverse-$Q$ factor on (B.14.1)
removes one of its two coefficient Tate
factors. Thus Frobenius acts on
$Z/(T-1)Z$ by $Q$ times the permutation
of the diagonal-$J_0$ orbit set in
$X\times X$, restricted to the antisymmetric
summand. This group has no invariant vector.
On each finite invariant partition,
a cycle of length $h$ would force
$(Q^h-1)v=0$ in a free $\mathbb Z_3$-module,
which implies $v=0$.
The inverse-limit orbit-module comparison
then gives the same vanishing integrally.
There is no extra orbit reversal: if a
translation in $J_0$ sends $(x,y)$ to
$(y,x)$, its square fixes $x$; freeness
and the absence of order-two elements in
$J_0$ imply $x=y$.
Together with $Z^{J_0}=0$, this proves
\[
                      H^1_{\rm cts}(G_k,Z)=0.              \tag{B.14.4}
\]

The difference between two lifts of the
fixed $B$-section is a continuous
$\mathscr I$-cocycle, for the reference
action just used.
By (B.14.3) its image in $M$ is a
coboundary. Lift a correcting element
from $M$ to $\mathscr I$ and conjugate
the section. Its cocycle now has values
in the central subgroup $Z$.
Equation (B.14.4) removes this last
cocycle by a central conjugation.
This proves uniqueness of the actual
cusp lift in the ramified projected case.

If a class-two section is unramified,
choose any lift of its geometric
Frobenius value to $\mathscr Q$.
Continuous powering defines an
unramified lift on the procyclic
unramified Galois quotient.
The ramified cases are exactly B.12
and the case just proved, giving the
full image assertion.
For a ramified unmarked cusp class,
its projected $B$-section is the unique
cusp class determined by its class-two
image: the nonzero inertia line identifies
the cusp, and the completed local
Kummer parameter is already injective
in the abelian cusp packet.
The uniqueness above therefore gives
the singleton fiber.
The marked fiber is precisely B.12--B.13.
The class-two tame identity, the correcting-conjugation
sign, and the finite orbit-sum certificates for the
nonzero marked central fiber are checked in
\texttt{code/bc\_marked\_checks.py} and its output.
\end{proof}

\begin{statement}{B.15}{Proved}{An explicit cleaned Heisenberg detector at seven}
Put $k=\mathbb F_7(\theta)$, $\theta^3=2$, and write
$[a,b,c]=a+b\theta+c\theta^2$.  On
$E_k:y^2+y=x^3$, let $\zeta=2$ and
\[
 \begin{split}
 P&=(3,2),\qquad P^-=(3,4),\\
 Q&=([4,1,1],[2,0,2]),\\
 S&=([0,6,5],[1,4,3]),\\
 T_1&=([3,3,3],[6,2,2]),\qquad T_2=-T_1 .
 \end{split}                                               \tag{B.15.1}
\]
The following functions are explicit:
\[
 \begin{aligned}
 a&=\frac{x^2+[1,6,5]x+[0,1,4]+[4,4,4]y}{x-3},&
 b&=x-3,\\
 f_1&=\frac{x^2+[2,3,5]x+[6,4,6]+[6,1,6]y}
                  {x-[3,3,3]},\\
 f_2&=\frac{x^2+[2,3,5]x+[0,3,0]+[1,6,1]y}
                  {x-[3,3,3]},\\
 m&=\frac{x^2+[2,5,2]x+[6,2,2]+[2,4,5]y}
                  {x-[3,3,3]},&
 f&=\frac{m}{f_1f_2^2(x-[3,3,3])}.
 \end{aligned}                                             \tag{B.15.2}
\]
Define $p(v)=\sum_{i=0}^5p_iv^i$ and
$h(v)=\sum_{i=0}^3h_iv^i$ by
\[
 \begin{array}{c|rrrrrr}
 i&0&1&2&3&4&5\\ \hline
 p_i&[3,0,0]&[0,6,6]&[6,1,5]&[4,0,2]&[4,1,2]&[1,5,4]\\
 h_i&[6,1,4]&[6,0,1]&[2,0,0]&[1,5,4]&&
 \end{array}                                               \tag{B.15.3}
\]
In $k(E)(v)$ with $v^3=b$, put
\[
 u(v)=\frac{p(v)+y}{v h(v)},\qquad
 r=f\,u(v)u(\zeta v)^2.                                    \tag{B.15.4}
\]
Then
\[
              v^3=b,\qquad w^3=a,\qquad z^3=r               \tag{B.15.5}
\]
defines a connected Galois cover of degree $27$, with
Heisenberg group of exponent three.  It is finite \'etale
over $U_k=E_k\setminus E_k[3]$.
The actual deck generators are
\[
 \begin{array}{c|ccc}
       &v&w&z\\ \hline
 \sigma&\zeta v&w&u(\zeta^2v)z/w\\
 \tau&v&\zeta w&z\\
 \rho&v&w&\zeta z ,
 \end{array}
 \qquad [\sigma,\tau]=\rho.                                 \tag{B.15.6}
\]
For the Fermat tangent parameter $s$ at $O$, the first
two radicals retain the leading coefficient $9$, whose
reduction is $2$.  The cleaned central radical satisfies
$r\sim s^4$.  Thus the based Puiseux choices are
\[
 v\sim\theta s^{-2/3},\qquad
 w\sim\theta s^{-2/3},\qquad
 z\sim s^{4/3}.                                            \tag{B.15.7}
\]
There is also an explicit Frobenius transport on the
compositum of the three coefficient-conjugate covers.
Let a superscript $[j]$ mean raising every coefficient
to $7^j$, with $j$ read modulo three, and use the same
choices (B.15.7) for $w_j,z_j$.  The arithmetic
Frobenius at the framed tangent is
\[
 \begin{split}
 \Phi(v)&=\zeta^2v,\qquad
 \Phi(w_j)=\zeta^2w_{j+1},\\
 \Phi(z_j)&=
       \frac{\zeta w_{j+1}z_{j+1}}
                    {u^{[j+1]}(\zeta v)},\qquad
 \Phi^3=1 .
 \end{split}                                               \tag{B.15.15}
\]
\end{statement}
\begin{proof}
All coefficients in this proof are computed in the
displayed finite field.  The group law gives
$3Q=P$ and $3S=T_2$.  The function $a$ is the normalized
Miller function
\[
                  \operatorname{div}(a)=3(Q)-(P)-2(O).
                                                               \tag{B.15.8}
\]
For example, multiplying the tangent line at $Q$ and the
line through $2Q,Q$, and dividing by the vertical lines
at $2Q,P$, gives (B.15.2); the common factor
$x-x(2Q)$ cancels in both numerator components.
This is an equality of rational functions with the
retained tangent leading coefficient, rather than an
identification of divisor classes.

The full norm identity is
\[
                    N_{k(E)(v)/k(E)}u=a.                  \tag{B.15.9}
\]
Here is a direct finite calculation of it.  Write
\[
 R(B)=[6,1,0]+[0,3,1]B+[0,2,2]B^2+[5,4,0]B^3.
\]
Expansion using $y^2+y=x^3$, $v^3=x-3$ gives
\[
 N(p(v)+y)=a(x-3)R(x-3),\qquad N(h(v))=R(x-3).
                                                               \tag{B.15.10}
\]
Since $N(v)=x-3$, (B.15.9) follows without a scalar
ambiguity.  The norm calculation can also be checked by
the cubic identity
$N(A+Bv+Cv^2)=A^3+bB^3+b^2C^3-3bABC$.

We next exhibit the removal of every auxiliary
ramification point.  The cubic $h$ is irreducible over
$k$.  Let $\alpha$ be its root and let
\[
 R_i=(\alpha^3+3,-p(\zeta^i\alpha)),\qquad i=1,2.
                                                               \tag{B.15.11}
\]
These are points of $E(k(\alpha))$; the candidate with
$i=0$ does not lie on $E$.  Their Frobenius orbits have
degree three.  Denote the associated closed-point
divisors by $\mathcal R_1,\mathcal R_2$.
Exact addition of the three conjugates gives
$\operatorname{Tr}(R_i)=T_i$.
The two successive line constructions for
$R_i+R_i^{343}+R_i^{343^2}=T_i$ give precisely $f_i$:
\[
 \operatorname{div}(f_i)=\mathcal R_i-(T_i)-2(O).
                                                               \tag{B.15.12}
\]
The coefficients descend to $k$ after cancellation of
the intermediate vertical factor; they are the
coefficients displayed in (B.15.2).
Likewise $m$ is the normalized Miller function with
divisor $3(S)-(T_2)-2(O)$.  Since $T_1=-T_2$, the
vertical function has divisor
$(T_1)+(T_2)-2(O)$.  Therefore
\[
 \operatorname{div}(f)
       =3(S)+6(O)-\mathcal R_1-2\mathcal R_2.               \tag{B.15.13}
\]

At $\mathcal R_i$, the valuations modulo three of
$u(v)u(\zeta v)^2$ are $i$, on all three sheets.
For example at $v=\alpha$ the denominator contributes
$-1$, and the numerator contributes $1$, $2$, or $0$
according to whether its zero occurs at
$\alpha,\zeta\alpha$, or $\zeta^2\alpha$.
The two actual points have indices $1,2$, giving
residues $1,2$.  Formula (B.15.13) cancels them.
At the three lifts of $Q$ the numerator orders are
$0,0,3$; the
central valuations are therefore multiples of three.
The extra zero $S$ in (B.15.13) also has order three.
Equations (B.15.8)--(B.15.10) list all other zeros and
poles: they lie over $P,P^-,O$.  The first two cubic
extensions are unramified at all interior points, and
the last radical now has valuation divisible by three
there.  Since the residue characteristic is seven,
the normalized cover is \'etale over $U_k$.

For completeness, the field degrees and deck morphisms
can be verified directly.  The classes of $a,b$ in
$k(E)^\times/k(E)^{\times3}$ are independent: the
valuation at $P^-$ first forces the coefficient of $b$
to vanish, and the valuation at $P$ then forces the
coefficient of $a$ to vanish.  Hence $[k(E)(v,w):k(E)]=9$.
Writing $\sigma(v)=\zeta v$, (B.15.9) gives
\[
             \frac{\sigma(r)}r
                   =\frac{u(\zeta^2v)^3}{a}.              \tag{B.15.14}
\]
If $r$ became a cube in $k(E)(v,w)$, the Kummer kernel
would give $r=a^j c^3$ in $k(E)(v)$ for some $j$.
Applying $\sigma$ and (B.15.14) would make
$a=(u(\zeta^2v)c/\sigma(c))^3$ in $k(E)(v)$, contradicting
the independence just proved.  The last degree is
therefore three.
Equations (B.15.14) and (B.15.9) prove that (B.15.6)
preserves all defining equations, that $\sigma^3=1$,
and that $[\sigma,\tau]=\rho$.  These are actual
automorphisms of the degree-$27$ field.

Finally use $t=-x/y\sim s/3$ at $O$.
The functions $a,b,f_1,f_2,m$ and the vertical factor
all have leading term $9s^{-2}$.  Thus
$u(v)\sim v$ and $f\sim s^6/3^6$.
The original central radical has leading term
$\zeta^2\,9s^{-2}$, so after the twist its coefficient
is $\zeta^2\,9/3^6=1$ in $k$ and its exponent is four.
This proves (B.15.7).  In particular no replacement of
the minimal-model tangent constant by $36$ is made.

To prove the last assertion, embed the fields in the
Puiseux field using (B.15.7).  Coefficientwise seventh
power, fixing $s^{1/3}$, acts on $k((s^{1/3}))$ with
order three, and fixes the Laurent expansions of
$x,y$.  It sends the chosen leading coefficient
$\theta$ to $\theta^7=\zeta^2\theta$.
Moreover (B.15.9) gives the exact identity
\[
 \frac{r^{[j]}(\zeta^2v)}{r^{[j]}(v)}
       =\frac{a^{[j]}}{u^{[j]}(\zeta v)^3}.
\]
Thus the proposed image of $z_j$ in (B.15.15) has the
correct cube.  Its leading coefficient is one, because
$w_{j+1}/u^{[j+1]}(\zeta v)$ has leading coefficient
$\zeta^{-1}$.  The uniqueness of the chosen cube root
with this leading coefficient proves the formula.
It respects every relation among the conjugate covers,
since it is the restriction of the actual coefficient
automorphism of the ambient Puiseux field.  Iterating
three times gives $w_j^3/N(u^{[j]})=1$ in the multiplier
of $z_j$, verifying $\Phi^3=1$ directly as well.

The exact polynomial identities, the finite-field
extensions, the three-conjugate additions, and the
displayed coefficients are independently replayed by
\texttt{code/bc\_miller\_norm\_rr.py} and
\texttt{code/bc\_miller\_heisenberg.py}, with their saved
JSON and text outputs.  This explicit based cover is a
detector for the remaining central-action calculation;
its construction and the conjugate-cover morphism
(B.15.15) do not by themselves assign its central
monodromy values on an uncomputed fixed Schreier basis.
\end{proof}

\begin{statement}{B.16}{Proved}{The integral depth-one map through the cyclic-nine cover}
Let $F$ be the free pro-three group on the standard
tripod inertia generators $x,y$, with
$[x,y]=xyx^{-1}y^{-1}$.
For $f\in F'$ write its metabelian image as
$[x,y]^{h(X,Y)}$, where $X=x-1$, $Y=y-1$ act by
conjugation.  There is an actual continuous
metabelian quotient in which
\[
 \begin{gathered}
 x\longmapsto(1,0),\qquad y\longmapsto(0,1),\\
 (m,a)(n,b)=(m+t^a n,a+b),\\
 m,n\in\mathbb Z_3[[t-1]],\qquad a,b\in\mathbb Z_3.
 \end{gathered}
                                                               \tag{B.16.1}
\]
In this quotient the image of $f$ is $(m_f,0)$, with
\[
                   m_f(t)=(1-t)h(0,t-1).                  \tag{B.16.2}
\]

Let $H_9=\ker(F\to\mathbb Z/9)$, where $x\mapsto0$,
$y\mapsto1$.  Its actual cover is
\[
 z^9=1-t_{\rm base},\qquad
 Y_9=\mathbb P^1\setminus\{0,\infty,\mu_9\}.
                                                               \tag{B.16.3}
\]
In $H_9^{\rm ab}/3$, write the image of $f$ as
$\sum_{j=0}^8c_j\,y^jxy^{-j}$; its $y^9$ coordinate
is zero.  Then
\[
 \begin{pmatrix}h_{0,3}\\h_{0,5}\end{pmatrix}
 =
 \begin{pmatrix}
 0&0&0&0&2&1&0&1&2\\
 0&0&0&0&0&0&2&2&2
 \end{pmatrix}
 \begin{pmatrix}c_0\\c_1\\c_2\\c_3\\c_4\\c_5\\c_6\\c_7\\c_8\end{pmatrix}
 \quad\text{in }\mathbb F_3.                              \tag{B.16.4}
\]

For the standard path $\beta$ from $\overrightarrow{01}$
to $\overrightarrow{10}$ and the convention
$f_g=\beta(g\beta)^{-1}$, these coordinates are explicit
framed Kummer coordinates on (B.16.3).
The relevant endpoint-to-start ratios are
\[
 r_0=9,\qquad r_j=\frac{-\zeta_9^j}{1-\zeta_9^j}
       \quad(1\leq j\leq8),\qquad \prod_{j=0}^8r_j=1.
                                                               \tag{B.16.5}
\]
Choose their compatible roots by continuing the units
$z-\zeta_9^j$ along the lifted path from $z=1$ to $z=0$,
using the actual start velocity $-1/9$ for $j=0$.
If
\[
 \frac{g(r_{\chi(g)^{-1}j}^{\,1/3^n})}
      {r_j^{\,1/3^n}}
     =\zeta_{3^n}^{\,\kappa_j(g)},                          \tag{B.16.6}
\]
where the index is taken modulo nine, then
$c_j=-\kappa_j(g)\pmod3$.
In particular (B.16.4)--(B.16.6) evaluate the two
coefficients by finite cyclic-cover Kummer data for
the full cyclotomic action; no restriction to
$\ker\chi$ is imposed.
\end{statement}
\begin{proof}
The group in (B.16.1) is a compact pro-three group.
The universal property of $F$ gives the stated
continuous homomorphism.  Its target is metabelian,
and conjugation by $x$ acts trivially on its additive
normal subgroup, while conjugation by $y$ multiplies
that subgroup by $t$.  Direct multiplication sends
$[x,y]$ to $(1-t,0)$.  Therefore the induced
commutator-module map sends $X$ to zero and $Y$ to
$t-1$, proving (B.16.2).

The lift of a positive loop about $t_{\rm base}=1$
in (B.16.3) sends $z$ to $\zeta_9z$.
Consequently the Schreier generators
$y^jxy^{-j}$ are the actual positive inertia loops
at $z=\zeta_9^j$, and $y^9$ is the positive inertia
loop at $z=0$.  The functions $z-\zeta_9^k$ and $z$
have valuations respectively $\delta_{jk}$ and zero
on the first nine generators, and respectively zero
and one on the last generator.  Their Kummer
characters thus give the actual dual coordinate
maps on $H_9^{\rm ab}/3$.
The Schreier basis has ten members, so these
coordinate maps are an isomorphism with
$\mathbb F_3^{10}$.
Since $f\in F'$, its total $y$ exponent is zero;
hence its $y^9$ coordinate is zero.

Reducing (B.16.2) modulo three and $t^9-1$ gives
\[
       \sum_{j=0}^8c_j(1+Y)^j=-Yh(0,Y)\pmod{Y^9},
                                                               \tag{B.16.7}
\]
because $(1+Y)^9-1=Y^9$ in characteristic three.
The coefficients of $Y^4,Y^6$ give
$h_{0,3}=-\sum_j\binom j4c_j$ and
$h_{0,5}=-\sum_j\binom j6c_j$, which are exactly
(B.16.4).
Only these indicated coefficients are recovered:
the $Y^8$ coefficient of $h$ is lost under
multiplication by $Y$ modulo $Y^9$.

The path $\beta$ lifts to the real interval from
$z=1$ to $z=0$.  At the start,
$z-1=-t_{\rm base}/9+\cdots$, and at the end
$z-1=-1$.  Their framed ratio is therefore nine.
For $j\ne0$ both endpoint values are ordinary
nonzero values, giving (B.16.5).
The product identity follows from
$\prod_{j=1}^8(1-\zeta_9^j)=9$ and
$\prod_{j=1}^8(-\zeta_9^j)=1$.
The continued roots obey the product identity as
well: their arguments occur in opposite pairs,
and the $j=0$ root is the positive real root of nine.

To check the sign in (B.16.6), choose one root of
the unit at the start and continue it along $\beta$
to the end.  The path $g\beta$ transports the
Galois conjugate start root to the Galois conjugate
end root, with the unit index permuted by $\chi(g)$.
Following $\beta$ and then $(g\beta)^{-1}$ multiplies
the start root by the inverse of the
endpoint-to-start Kummer multiplier.  Thus the
positive inertia coordinate of $f_g$ is
$-\kappa_j(g)$, as asserted.
This argument constructs the maps for every $g$;
it does not extend a formula stated only on
$\ker\chi$ by an unspecified normalization.

The program \texttt{code/bc\_depth\_one.py} checks
the actual Schreier coordinates against the
$X$ coefficient of the abelianized Fox derivative
$D_y(f)=Xh(X,Y)$ on explicit commutator words.
It also verifies the exact cyclotomic-field
identities (B.16.5), their degree-nine to
degree-three transition products, the Frobenius-seven
index permutation, and the two rows in (B.16.4).
The saved output has status \texttt{PASS}.
\end{proof}

\begin{statement}{B.17}{Proved}{The actual norm-289 depth-one Frobenius coordinates}
Let $g_{17}$ be an arithmetic Frobenius of norm $289$
at the prime above seventeen of $K_0=\mathbb Q(\zeta_3)$,
viewed through the specified tangential arithmetic action.
Use the path and commutator conventions of B.16, and write
the metabelian Ihara word as
$f_{g_{17}}=[x,y]^{h(X,Y)}$.
Then
\[
 h(0,Y)\equiv-Y^6-Y^7\pmod{(3,Y^8)},
 \qquad h_{0,3}=h_{0,5}=0\quad\text{in }\mathbb F_3.
                                                               \tag{B.17.1}
\]
An explicit residue-field embedding is
\[
 \begin{gathered}
 k=\mathbb F_{17}[\theta]/(\theta^2-3),\\
 \overline\zeta_9=7+4\theta,\qquad
 \overline\zeta_3=\overline\zeta_9^3=8+15\theta.
 \end{gathered}
                                                               \tag{B.17.2}
\]
For the nine framed ratios of B.16 the actual local
Kummer coordinates and geometric cusp coordinates are
\[
 \begin{aligned}
 (\kappa_j(g_{17}))_{j=0}^8
       &=(0,1,2,0,1,2,0,1,2),\\
 (c_j)_{j=0}^8&=(0,2,1,0,2,1,0,2,1).
 \end{aligned}
                                                               \tag{B.17.3}
\]
The entries are in $\mathbb F_3$. These values hold for
each prime of $\mathbb Q(\zeta_9)$ above seventeen,
with its compatible embedding of $\zeta_3=\zeta_9^3$.
\end{statement}
\begin{proof}
The integer three is a nonsquare modulo seventeen, and
direct multiplication in (B.17.2) gives
$\overline\zeta_9^9=1$,
$\overline\zeta_9^3\ne1$, and
$\overline\zeta_9^{17}=\overline\zeta_9^{-1}$.
Choose the embedding of $\mathbb Q(\zeta_9)$ into
$\overline{\mathbb Q}_{17}$ that reduces to (B.17.2).
Its completion is the unramified quadratic extension
of $\mathbb Q_{17}$, which is also the completion of
$K_0$. Every $r_j$ in (B.16.5) is a unit there.
The polynomial $Z^3-r_j$ is finite etale on the local
unit locus: its derivative is a unit at every root.
Its roots therefore lie in finite unramified extensions.

Let $\rho_j$ be the cube root selected by the actual
path of B.16. Reduction is injective on $\mu_3$.
Arithmetic Frobenius acts on its residue as the
$289$th power, so the exact local comparison is
\[
 \overline{\frac{g_{17}(\rho_j)}{\rho_j}}
    =\overline\rho_j^{288}
    =\overline r_j^{96}
    =\overline\zeta_3^{\,\kappa_j(g_{17})}.
                                                               \tag{B.17.4}
\]
Here $g_{17}$ fixes $\mu_9$, so there is no index
permutation in (B.16.6). It also fixes $\mu_3$;
changing a cube root by an element of $\mu_3$ does
not change (B.17.4). Thus this calculation computes
the prescribed based cocycle itself modulo three.
Changing the representative by local inertia also
leaves it unchanged because these unit extensions
are unramified.

We can evaluate every symbol without choosing any
radical extension. For $1\leq j\leq8$, the
seventeen-power Frobenius sends $\overline r_j$ to
$\overline r_{-j}$, with indices modulo nine, and
\[
 \overline r_j\overline r_{-j}\in\mathbb F_{17}^{\times},
 \qquad
 \frac{\overline r_j}{\overline r_{-j}}
     =-\overline\zeta_9^j.
                                                               \tag{B.17.5}
\]
Every element of $\mathbb F_{17}^{\times}$ has
trivial $96$th power. Since $96\equiv6\pmod9$,
(B.17.5) gives
\[
 \kappa_j+\kappa_{-j}=0,
 \qquad \kappa_j-\kappa_{-j}=2j
       \quad\text{in }\mathbb F_3.
\]
Thus $\kappa_j=j\pmod3$. The ratio $r_0=9$
also has trivial symbol, which proves the first
vector in (B.17.3). The actual sign comparison of
B.16 gives $c_j=-\kappa_j$ and proves the second.
The argument uses only primitivity of the chosen
ninth root and its cube, so it applies to every
embedding above seventeen.

For completeness, the entire specialization through
this precision follows from the same nine symbols.
In $\mathbb F_3[t]/(t^9-1)$,
\[
 \sum_{j=0}^8c_jt^j
   =-t\frac{d}{dt}\left(\sum_{j=0}^8t^j\right)
   =t(t-1)^7.
\]
With $t=1+Y$, equation (B.16.7) becomes
$-Yh(0,Y)=Y^7+Y^8\pmod{Y^9}$.
Division by $Y$ at the stated precision proves
(B.17.1). No coefficient of degree eight or higher
is assigned by this division.

The standalone program
\texttt{code/bc\_depth\_one\_frob17.py} executes
all field operations in (B.17.2), evaluates the nine
symbols by $96$th powers, verifies (B.17.5),
reconstructs the cyclic group-ring polynomial and
the two rows of B.16, and repeats the computation
for all six primitive ninth roots in $k$.
The saved JSON and text outputs have status
\texttt{PASS}.
\end{proof}

\begin{statement}{B.18}{Proved}{The full finite arithmetic image and the exact Frobenius action at seventeen}
Use the actual source basis of A.38 and the
$c$-adapted basis $m,n$ of A.31. Let
\[
 \rho:P_3\longrightarrow
 \operatorname{Aut}\bigl(N/(N^3\gamma_3N)\bigr)
\]
be the actual arithmetic action. Its image $G$
has order $243$, nilpotency class three, and
exponent nine. In the $10+45$ coordinates of A.38,
its linear image has order nine, and its
identity-linear kernel has dimension three over
$\mathbb F_3$. The exact dimensions are
\[
 \begin{array}{c|ccc}
   &\gamma_2(G)&\gamma_3(G)&\gamma_4(G)\\\hline
 \dim_{\mathbb F_3}&2&1&0
 \end{array}
 \qquad
 \begin{array}{c|ccc}
   &D_2(G)&D_3(G)&D_4(G)\\\hline
 \dim_{\mathbb F_3}&3&2&0.
 \end{array}
                                                               \tag{B.18.1}
\]
In particular $D_4(P_3)$ acts trivially on this
specified geometric quotient. Put
$M=\rho(m)$ and $N_*=\rho(n)$. Then the entire
actual Frobenius action at seventeen is
\[
 \rho(\varphi_{17})
   =M^6[N_*,[N_*,M]]=M^6\ne1.
                                                               \tag{B.18.2}
\]
It has order three. Its Johnson block has rank two
and annihilates all nine cusp vectors. The full
matrix is retained in
\texttt{bc\_ihara\_finite\_image\_certificate.npz}.
\end{statement}
\begin{proof}
An automorphism in these coordinates has matrix
\[
 \begin{pmatrix}A&0\\J&\bigwedge^2 A\end{pmatrix}.
\]
Its identity-linear subgroup is the additive group
$\operatorname{Hom}(\mathbb F_3^{10},
\bigwedge^2\mathbb F_3^{10})$.
For such a matrix $X$, conjugation is the exact map
\[
                  X\longmapsto
            (\bigwedge^2 A)XA^{-1}.                       \tag{B.18.3}
\]
The script below reconstructs $\rho(u)$ and $\rho(c)$
from the A.38 word and its literal Schreier map,
forms $\rho(v)=\rho(c)\rho(u)\rho(c)$, and applies
the actual square-root source maps defining $m,n$.
These reconstructed matrices agree with the
independently saved arithmetic certificate.

The two linear matrices of $M,N_*$ commute, each
has order three, and their nine products are
distinct. Define $Z_0,Z_1,Z_2$ to be the additive
Johnson blocks, respectively, of
\[
 M^3,\qquad [M,N_*],\qquad [M,[M,N_*]].
\]
Exact substitution gives three linearly independent
vectors and the conjugation table
\[
 \begin{array}{c|ccc}
       &Z_0&Z_1&Z_2\\\hline
 \operatorname{Ad}(M)&Z_0&Z_1+Z_2&Z_2\\
 \operatorname{Ad}(N_*)&Z_0&Z_1&Z_2.
 \end{array}                                                \tag{B.18.4}
\]
It also gives $N_*^3=1$. The subgroup generated
by the three $Z_i$ is therefore normal, is
elementary abelian of order $27$, and is contained
in the actual image. After quotienting by it,
the images of $M,N_*$ commute and have order
three. The linear image has order nine, so the
quotient is exactly $(\mathbb Z/3)^2$. This proves
$|G|=243$, and shows that the displayed subgroup
is the whole identity-linear kernel.

The normal closure of $[M,N_*]$ is exactly
$\langle Z_1,Z_2\rangle$ by (B.18.4). Its next
commutator is $\langle Z_2\rangle$, which is
central. Thus the lower-central assertions in
(B.18.1) follow as actual group statements.
Every cube lies in the identity-linear kernel,
and that kernel has exponent three; consequently
every element has ninth power one. Since $M^3$
is nontrivial, the exponent is nine. The cubes
span $\langle Z_0,Z_2\rangle$, as verified on
the $243$ distinct normal forms. This also proves
the two asserted Zassenhaus dimensions.

The defining formula for this filtration gives
\[
 D_4(P_3)=
 \overline{\gamma_4(P_3)(\gamma_2P_3)^3P_3^9}.
\]
Each factor has trivial image by the preceding
calculations. Equation (A.31.3), transported by
the actual automorphism in A.38, therefore gives
the equality in (B.18.2), not only a leading-term
comparison. Table (B.18.4) gives
$[N_*,[N_*,M]]=1$. The matrix $M^6$ has order
three, and its $45\times10$ Johnson block has
rank two. Multiplication by the nine actual
cusp columns of the Schreier certificate gives
zero in every case.

The standalone
\texttt{code/bc\_ihara\_finite\_image.py}
computes the full normal commutator module by
the exact maps (B.18.3), checks (B.18.4), and
independently enumerates the subgroup generated
by $M,N_*$ using literal matrix multiplication.
Both computations give the same $243$ elements.
It verifies all ninth powers, the cube span,
the Frobenius word, its rank and cusp values,
and saves every indicated basis and matrix.
Its JSON and text outputs have status
\texttt{PASS}.
\end{proof}

\begin{statement}{B.19}{Proved}{A finite precision determining the integral local kernel}
Let $M_2=\bigwedge^2V$, and let
\[
 \mathcal D=M_2^c\oplus
       \ker(1+\rho(n)c:M_2\longrightarrow M_2)
\]
be the rank-$45$ lattice of normalized source cocycles
used in A.39. Define the actual Frobenius evaluation
\[
 E_\varphi:\mathcal D\longrightarrow M_2^c,
 \qquad E_\varphi(\gamma)=\gamma(\varphi_{17}),
 \qquad R=\ker E_\varphi.
\]
The integral operator
$B=(F_\varphi-1)|M_2^c$ has three-adic Smith
exponents
\[
 \underbrace{1,\ldots,1}_{6},\quad
 \underbrace{2,\ldots,2}_{7},\quad
 \underbrace{3,3}_{2},\quad
 \underbrace{4,\ldots,4}_{6}.
                                                               \tag{B.19.1}
\]
Identify $R/3R$ with its image in $\mathcal D/3\mathcal D$.
Then its exact finite-precision description is
\[
 \frac{R}{3R}
 =\operatorname{im}\left(
 \ker\left(
 \mathcal D/243\mathcal D
       \xrightarrow{\ E_\varphi\ }
 M_2^c/243M_2^c\right)
 \longrightarrow\mathcal D/3\mathcal D\right).
                                                               \tag{B.19.2}
\]
More explicitly, for every integral $d\in\mathcal D$
with $E_\varphi(d)\in243M_2^c$, put
\[
 h=B^{-1}E_\varphi(d)\in3M_2^c,
 \qquad \gamma=d-\delta h,
 \qquad (\delta h)(g)=g(h)-h.
                                                               \tag{B.19.3}
\]
Then $\gamma\in R$ and $\gamma\equiv d\pmod{3\mathcal D}$.
The matrix $B$ and the bound (B.19.1) are computed
in the attached certificate. Evaluating the separate
arithmetic matrix $E_\varphi\bmod243$ remains the
finite calculation needed to list the resulting
$24$-dimensional subspace individually.
\end{statement}
\begin{proof}
At seventeen, cubing permutes $\mathbb F_{17}$.
Consequently the reduction $y^2+y=x^3-7$ has one
$x$ for each $y$, and its point count, including
infinity, is eighteen. Its Frobenius polynomial
is $X^2+17$, so the norm-$289$ Frobenius on $T$
is exactly $-17$. All cusps are rational over
$K_0$, and its action on $W$ is multiplication
by $289$. The actual graph-lattice comparison
(A.7.4) therefore gives
\[
 F_\varphi|V=
 \begin{pmatrix}
 289I_8&102\Psi\\0&-17I_2
 \end{pmatrix}.
                                                               \tag{B.19.4}
\]
We take its full exterior square and restrict
to the $c$-fixed lattice. The latter has rank
$21$. The projection to $\bigwedge^2T$ is zero
on it, because $c$ acts there by $-1$; hence
it is also the fixed lattice inside $\mathcal K$.

The program constructs a basis from $21$ independent
columns modulo three of $1+c$ on $M_2$. This is
an actual $\mathbb Z_3$ basis of the fixed lattice:
$\frac12(1+c)$ is an integral idempotent, and
the selected columns span its reduction modulo
three, so Nakayama's lemma gives the full lattice.
An invertible minor of three-adic unit determinant
gives the coordinate restriction of (B.19.4).
Both the rectangular matrix in $M_2$ and its
$21\times21$ coordinate matrix have the exponents
(B.19.1). An independent elimination modulo
$3^8$ gives the same list. The determinant check is
\[
 \det B=83520^{13}\,4914^8,
 \qquad v_3(\det B)=50.
\]
Indeed the rational fixed subspace has thirteen
toric wedges, of eigenvalue $289^2$, and eight
mixed wedges, of eigenvalue $-17\cdot289$.
The inverse matrix verifies that $81B^{-1}$
is integral over $\mathbb Z_3$, while
$27B^{-1}$ is not. This proves the stated bound
with its actual integral comparison.

For $h\in M_2^c$, the coboundary $\delta h$
is normalized at $c$ and belongs to $\mathcal D$.
Its evaluation is exactly
\[
                  E_\varphi(\delta h)=Bh.
\]
Thus $\operatorname{im}E_\varphi$ contains
$B(M_2^c)$, which contains $81M_2^c$ by
(B.19.1). The formula (B.19.3) is therefore
integral and has $h\in3M_2^c$ whenever
$E_\varphi(d)$ is divisible by $243$.
It gives $E_\varphi(\gamma)=0$ and leaves
the reduction of $d$ modulo three unchanged.
Conversely every element of $R$ reduces to
an element of the displayed finite kernel.
Since the target of $E_\varphi$ is torsion-free,
$R\cap3\mathcal D=3R$. This proves both the
claimed identification and (B.19.2).

The exact reconstruction, point count, two Smith
computations and inverse-denominator check are
executed in
\texttt{code/bc\_frobenius17\_precision.py}.
Its JSON and text outputs have status
\texttt{PASS}, and include the fixed-lattice
basis and the full integral matrix $B$.
\end{proof}

\begin{statement}{B.20}{Proved}{An individually specified class-two target with no class-four lift}
Use the actual arithmetic basis $m,n,c$ of A.38 and
A.31. Write the ten integral A.7 basis vectors as
\[
 \begin{split}
 (v_0,\ldots,v_7)
   &=(e_P,e_{Q_c},e_A,e_{\bar A},e_B,e_{\bar B},e_C,e_{\bar C}),\\
 v_8&=(\Psi(1,0)/3,(1,0)),\qquad
 v_9=(\Psi(0,1)/3,(0,1))
       \quad\text{in }L_\Psi\subset W_{\mathbb Q_3}\oplus T_{\mathbb Q_3}.
 \end{split}                                                \tag{B.20.1}
\]
Let $C_2$ be the actual action of $c$ on
$M_2=\bigwedge^2V$. Define the integral vectors
\[
 b_0=3v_8\wedge v_9+
       \sum_{i=0}^7
       \bigl(\Psi_{i,2}v_i\wedge v_8
             -\Psi_{i,1}v_i\wedge v_9\bigr),
 \qquad b=(1-C_2)b_0.                                     \tag{B.20.2}
\]
Here the columns of $\Psi$ are numbered one and two.
In the basis $v_i\wedge v_j$, $i<j$, all nonzero
coordinates of $b$ are the following; every other
coordinate is zero:
\[
 \begin{array}{cr|cr|cr}
 (i,j)&b_{ij}&(i,j)&b_{ij}&(i,j)&b_{ij}\\\hline
 (0,8)&-2&(1,8)&2&(2,6)&-1\\
 (2,7)&1&(2,8)&2&(2,9)&2\\
 (3,6)&1&(3,7)&-1&(3,9)&-2\\
 (4,6)&1&(4,7)&-1&(4,8)&-2\\
 (4,9)&-2&(5,6)&-1&(5,7)&1\\
 (5,9)&2&(6,8)&2&(6,9)&-2\\
 (7,8)&-2&(7,9)&2&(8,9)&6.
 \end{array}                                                \tag{B.20.3}
\]
The prescriptions
\[
             \gamma(m)=0,\qquad \gamma(n)=b,
             \qquad \gamma(c)=0                           \tag{B.20.4}
\]
define an actual continuous $M_2$-cocycle unramified
outside three. Let $s_{O,17}$ be the actual $O$-cuspidal
section with tangent parameter seventeen. In
$A=N/\gamma_3N$, put
\[
                        s_b(g)=\gamma(g)s_{O,17}(g).
                                                               \tag{B.20.5}
\]
This is an actual class-two section, with abelian
class $\kappa_{17}e_O$. Its restriction at seventeen
equals the specified local cuspidal section exactly.
It has a global lift to $N/\gamma_4N$, but no global
lift to $N/\gamma_5N$. Consequently it has no full,
metabelian, or marked-depth lift.

The projection of (B.20.4) to $\bigwedge^2T$ has
the exact generator values $0,6$. If
$a=\kappa_3(n)\in\mathbb Z_3^\times$ in this actual
source basis, this cocycle is $(6/a)\kappa_3$.
Only $a\equiv1\pmod3$ is used; its integral value
has not been replaced by one.
\end{statement}
\begin{proof}
We first prove the integral local condition. On $V$,
the actions of $m,n$ are respectively $G_uG_v$ and
$G_vG_u^{-1}$ from (A.7.4)--(A.7.5). They commute,
and the first commutes with $c$. Set
\[
 L_0=\ker(n-1:M_2\to M_2)\cap\ker(c+1:M_2\to M_2).
\]
This is a closed saturated lattice stable under
$m,n,c$ and their inverses. Direct substitution
in the exact graph matrices gives
\[
 n(b_0)=b_0,\qquad n(b)=b,
 \qquad c(b)=-b.                                         \tag{B.20.6}
\]
For the first equality, the rational comparison
is also explicit: the vector $3$ times the wedge
of the two pure proper vectors has the mixed
coordinates in (B.20.2), proper coefficient three,
and toric coefficient
$\frac13\Psi_1\wedge\Psi_2$. Subtracting this
last toric vector gives exactly $b_0$.
Both pure proper determinant and toric wedges are
fixed by $n$, so the equality transfers through
the actual graph embedding. No reduction of a
rational splitting modulo three is taken.

The full lattice is
\[
           L_0=(\bigwedge^2W)^{c=-1}\oplus\mathbb Z_3b,
           \qquad \operatorname{rank}L_0=16.             \tag{B.20.7}
\]
Here is an integral verification. The toric
anti-invariant lattice has rank fifteen.
Choose the lexicographically first independent
columns modulo three of $1-c$ on its $28$
wedges. These are an actual integral basis by
the order-two projector and Nakayama's lemma.
Adjoining $b$ gives sixteen columns independent
modulo three. The stacked rational matrices
$n-1,c+1$ have rank $29$, so these columns span
their sixteen-dimensional common rational kernel.
Their unit minor makes their integral span
primitive in $M_2$, proving (B.20.7).
The code also computes the exact matrix of $m$
on this basis and its inverse; both are integral
over $\mathbb Z_3$. This verifies the stated
stability on the actual lattice.

Since $m,n$ freely generate $P_3$, their prescribed
values give a continuous crossed homomorphism.
The identities $c(m)=m$, $c(n)=n^{-1}$ and
(B.20.6) prove its exact $c$-equivariance, so
it extends with value zero at $c$. Every value
lies in $L_0$. The actual norm-$289$ Frobenius
$\varphi_{17}$ is fixed by $c$, hence
\[
 \gamma(\varphi_{17})=c\gamma(\varphi_{17})
                    =-\gamma(\varphi_{17})=0.
                                                               \tag{B.20.8}
\]
The final equality uses that two is invertible.
Thus (B.20.4) lies in the exact integral lattice
$R$ of A.39, not merely in a finite approximation
to it. Its inflation has zero tame value at
seventeen and zero value at the normalized local
involution as well. These three values prove the
exact local equality in (B.20.5).

We give a finite functional detecting the ensuing
gluing obstruction. Let $X_0,\ldots,X_9$ be the
ordered Schreier generators in (A.38.3), and use
the integral degree-three basis
\[
 H_{\mu,i}=[[X_i,X_{\min\mu}],X_j],
\]
where $\mu$ is a nondecreasing multiset of size
three, $i>\min\mu$ occurs in it, and $j$ is the
remaining index after removing $i$ and one
occurrence of $\min\mu$. Write $q_{\mu,i}$ for
the corresponding coordinate functional.
Over $\mathbb F_3$, define
\[
 \begin{split}
 \ell=2\bigl(&q_{001,1}+q_{002,2}+q_{003,3}+q_{004,4}\\
             &+q_{088,8}+q_{188,8}+q_{288,8}
                       +q_{388,8}+q_{488,8}\bigr).
 \end{split}                                                \tag{B.20.9}
\]
All these indices refer to the Schreier basis,
not to an unidentified change of proper lifts.
The program uses the exact unimodular
$A.7$-to-Schreier matrix $P$ and the actual
map $\bigwedge^2P$ on $b$. It verifies that this
map intertwines the full $m,n,c$ actions on
$M_2$ before evaluating the source cochain.

With the independent Magnus evaluator of
\texttt{bc\_gluing\_audit.py}, the functional
has the following exact values. It kills the
entire global $M_3$ restriction image, all
twenty-one retained local-conjugator parameters,
and all sixty-eight unknown-word contributions
on the cocycle (B.20.4). Its value on the known
Frobenius contribution is
\[
                           \ell(\Delta)=1.
                                                               \tag{B.20.10}
\]
Only three of those unknown maps are nonzero
on the general source, and all three vanish on
the whole lattice $L_0$ modulo three. The known
map vanishes on its fifteen-dimensional toric
summand and is nonzero on $b$. In the fixed
$143$-dimensional quotient of the independent
audit, its vector has $87$ nonzero entries.
All twenty-seven remaining arithmetic coefficient
possibilities give exactly this same vector.

These are evaluations of the actual global
gluing map of A.39. Changing the integral
class-three lift adds a global $M_3$ correction,
which (B.20.9) kills. The retained source
truncation and all local comparison parameters
are those independently verified there, so
(B.20.10) applies to the actual arithmetic
Frobenius. A.33 gives a global class-three lift;
A.40 proves that this nonzero gluing class
prevents every class-four lift. The actual
quotient maps and local cusp comparisons in
A.39--A.40 give the full, metabelian and marked
non-lifting conclusions on section classes.

Finally the proper projection of $b$ is six
by (B.20.3). The positive-root Kummer cocycle
$\kappa_3$ has value zero at $c$, and hence at
the $c$-fixed generator $m$. Its value at $n$
is a unit by (A.31.6). Equality on the actual
free source generators proves the asserted
identity with $(6/a)\kappa_3$.

The standalone
\texttt{code/bc\_explicit17.py} executes the
integral lattice, saturation, action and basis
comparisons, then every stated functional and
gluing evaluation. It saves all $45$ coordinates,
the exact rank-$16$ basis, the transition matrix,
the full nonzero vector and the nine-term
functional. Its JSON, NPZ and text certificates
have status \texttt{PASS}.
\end{proof}

\begin{statement}{B.6}{Proved}{The complete global one-prime image and its fibers}
Fix a rational prime $q\ne3$, put
$A=N/\gamma_3N$, and let $s$ be a continuous
class-two section unramified outside
$S=\{3,q,\infty\}$. Choose its exact cocycle
representative with $s(c)=1$, and write $\alpha$
for its projection to $V$. Define
\[
 \begin{split}
 H_0&=N,\qquad H_{n+1}=H_n^3[H_n,H_n],\\
 P_n&=N/H_n,\qquad A_n=P_n/\gamma_3P_n,\qquad
 J_n=\gamma_3P_n,\\
 S_n&=\{3,\infty\}\cup
 \{q:\alpha|_{I_q}\bmod3^n\ne0\}.
 \end{split}                                               \tag{B.6.1}
\]
All subgroups are closed. The exact finite target
$s_n$ is the image of $s$ in $A_n$.

Use the explicitly constructed finite based action on
$P_n$ and its splitting field $F_n$ from A.43.
Let $K_n$ be the compositum of $F_n$ with the
field of the target graph $(s_n,\rho_{A_n})$.
If $s_n$ is ramified outside $S_n$, the test fails.
Otherwise form the finite Galois field $\Omega_n$
as the compositum of $\Omega_{n-1}$ and all Galois
extensions $E/\mathbb Q$ containing $K_n$,
unramified outside $S_n$, of degree at most
$|P_n|[F_n:\mathbb Q]$, starting with
$\Omega_0=\mathbb Q$.
The discriminant and polynomial bounds in
(A.43.4)--(A.43.10) give a terminating finite
enumeration of these fields.

Let $X_n$ be the finite set of functions
$z:\operatorname{Gal}(\Omega_n/\mathbb Q)\to P_n$
satisfying the actual cocycle equation,
$z(c)=1$, and $z\bmod J_n=s_n$ literally.
For every $k\le n$, require its projection to
$P_k$ to be trivial on inertia at the primes
of $S_n\setminus S_k$. These conditions give
the actual transition maps of A.43. Then
\[
 \begin{split}
 [s]\in\operatorname{im}\bigl(\operatorname{Sec}(N)
          \longrightarrow\operatorname{Sec}(A)\bigr)
 &\Longleftrightarrow X_n\ne\varnothing\quad\text{for all }n\\
 &\Longleftrightarrow
 [s]\text{ lifts to }N/\gamma_rN
                    \quad\text{for all }r\ge3.
 \end{split}                                               \tag{B.6.2}
\]
Every such full or nilpotent lift is unramified
outside $S$. The complete fiber on section classes is
\[
                 \varprojlim_n X_n/J_n^c,                  \tag{B.6.3}
\]
where $h\in J_n^c$ acts by
$z(g)\mapsto h^{-1}z(g)g(h)$.
Thus (B.6.2) is an exact sequence of terminating
finite arithmetic tests, and (B.6.3) retains all
compatible conjugacy data. A bounded stopping rule
for every positive instance is not asserted.

The necessary local condition is precisely the
unramified class or a member of the actual cusp
packets (B.8.1). It is sufficient at seven by
B.10, and on the prime sets $\Sigma$ and $\Xi$
by A.16 and A.25. It is insufficient at seventeen.
The integral section of B.20, with abelian class
$\kappa(17)e_O$ and central generator values
$\gamma(m)=0$, $\gamma(n)=b$, has exactly the
local $O$-cuspidal class-two restriction. It has
a global $N/\gamma_4N$ lift and no
$N/\gamma_5N$ lift, hence no full, metabelian,
or marked-depth lift.
\end{statement}
\begin{proof}
For the full group, the three quotient maps in
(A.43.2) specialize exactly to (B.6.1). The image
of $H_n$ in $V$ is $3^nV$, by induction on the
Frattini operation. Hence
$A_n^{\rm ab}=V/3^nV$, and the finite target table
determines the displayed set $S_n$.
The actual quotient map sends $\gamma_3N$ onto
$\gamma_3P_n$, which proves the asserted kernel
$J_n$. This identifies every group and action
in the finite test with its A.43 counterpart.

The inherited ramification conditions ensure that
projection to $P_k$ factors through $\Omega_k$;
they are retained even when the prime $q$ first
enters the sets $S_n$ at a later precision.
When $\alpha(I_q)=0$, it never enters, so an
inertia value visible only in the central
class-two target is rejected as soon as it is
detected in some $A_n$. This is consistent with
B.7's detection of full inertia in $V$.
A.43 now gives the first equivalence in (B.6.2)
and the exact fiber (B.6.3). Its comparison
uses the surjections
$\ker(N\to A)^c\twoheadrightarrow J_n^c$
and the trivial global geometric centralizer
of the target in $A$. It therefore identifies
global section classes, without replacing the
conjugation kernel by a finite-level centralizer.
B.11 gives the second equivalence and the common
ramification bound, including for separately
chosen nilpotent lifts.

For the obstruction coordinates requested in
Task B.4, entry into $A$ is the residue,
tame-relation and convergence criterion of
Part A. The next central coefficient is the
actual free lattice
$M_3=\gamma_3N/\gamma_4N$. Its residue complex
at $q$ is
\[
 [\,M_3(-1)\xrightarrow{F_q-1}M_3(-1)\,],
 \qquad F_qt=qt
 \text{ for a tame inertia value }t\in M_3.                \tag{B.6.4}
\]
Here $F_q$ in the complex acts on the twisted
module and the second $F_q$ acts on $M_3$;
the equations are the same under the chosen
tame generator. A.33 computes the continuous
obstruction using this localization map. For
the fixed finite set $S$, the finite-coefficient
first-cohomology groups are finite, so their
first derived limit is zero. Actual local
cuspidal lifts kill the local degree-two
obstructions and give a global class-three
lift. Its remaining obstruction to having
the prescribed local cuspidal classes is
exactly the restriction cokernel (A.33.2).
A.39 evaluates a nonzero family in that
cokernel; A.40 proves the class-four exclusion;
B.20 specifies an individual member with all
integral coordinates and an explicit nine-term
finite functional. These facts prove the
stated failure of local sufficiency.

The finite tests (B.6.2) retain every subsequent
geometric layer and global relation. They thus
give the full one-prime image at seventeen
and at every other prime, including the
classes for which the local condition alone
does not decide lifting. B.1, B.3 and B.9
separately give all three global cusp packets
and their inherited central origins, as
required in Task B.1.
\end{proof}

